% ================================================================
%  第 1 章  理论的起点
% ================================================================
\chapter{理论的起点}
\label{chap:1}

本章中, 我们首先把极小曲面方程作为图上面积泛函的 Euler--Lagrange 方程推导
出来. 随后我们导出极小曲面方程的参数形式(第一变分公式). 本章的重点是
\newterm{minimal-surface}{极小曲面}{minimal surface}的一些基本性质, 包括
面积的单调性公式与 Bernstein 定理. 我们也提到一些例子. 在倒数第二节中, 我们
导出第二变分公式、稳定性不等式, 并定义极小曲面的 Morse 指标. 在最后一节中,
我们引入多值极小图, 它在后面讨论 \cite{CM3}--\cite{CM7} 的结果时将起主要
作用. 我们还会给出一个来自 \cite{CM18} 的螺旋型极小曲面的局部例子(像螺旋面
那样), 它可以分解为多值图, 但螺旋的速率远不是常数.

% ================================================================
\section{极小曲面方程与极小子流形}
\label{sec:1.1}

\mtpara{图与极小曲面方程}
设 $u:\Omega\subset\R^2\to\R$ 是 $C^2$ 函数, 考虑函数 $u$ 的图
\begin{equation}
\Graph_u=\bigl\{(x,y,u(x,y))\mid (x,y)\in\Omega\bigr\}.
\label{eq:1.1}
\end{equation}
则面积为
\begin{align}
\Area(\Graph_u)
  &=\int_{\Omega}\bigl|(1,0,u_x)\times(0,1,u_y)\bigr| \notag\\
  &=\int_{\Omega}\sqrt{1+u_x^2+u_y^2}
   =\int_{\Omega}\sqrt{1+|\nabla u|^2},
\label{eq:1.2}
\end{align}
而(向上指的)单位法向量为
\begin{equation}
N=\frac{(1,0,u_x)\times(0,1,u_y)}
        {\bigl|(1,0,u_x)\times(0,1,u_y)\bigr|}
 =\frac{(-u_x,-u_y,1)}{\sqrt{1+|\nabla u|^2}}.
\label{eq:1.3}
\end{equation}
把它用到单参数图族 $\Graph_{u+t\eta}$ 上, 其中 $\eta|_{\partial\Omega}=0$,
$t$ 为参数, 便得到
\begin{equation}
\Area(\Graph_{u+t\eta})=\int_{\Omega}\sqrt{1+|\nabla u+t\nabla\eta|^2};
\label{eq:1.4}
\end{equation}
于是面积泛函在图上的 $u$ 处沿方向 $\eta$ 的方向导数为
\begin{align}
\frac{d}{dt}\Big|_{t=0}\Area(\Graph_{u+t\eta})
  &=\int_{\Omega}\frac{\langle\nabla u,\nabla\eta\rangle}
        {\sqrt{1+|\nabla u|^2}} \notag\\
  &=-\int_{\Omega}\eta\,\dive\Bigl(\frac{\nabla u}
        {\sqrt{1+|\nabla u|^2}}\Bigr).
\label{eq:1.5}
\end{align}
因此, $u$ 的图是面积泛函的临界点, 当且仅当 $u$ 满足散度形式的方程
\begin{equation}
\dive\Bigl(\frac{\nabla u}{\sqrt{1+|\nabla u|^2}}\Bigr)=0.
\label{eq:1.6}
\end{equation}
方程 \eqref{eq:1.6} 是\newterm{mse}{极小曲面方程}{minimal surface equation}
的散度形式, 它也可以写成
\begin{align}
0
 &=\bigl(1+|\nabla u|^2\bigr)^{\frac32}
   \Bigl[\Bigl(\frac{u_x}{\sqrt{1+|\nabla u|^2}}\Bigr)_x
        +\Bigl(\frac{u_y}{\sqrt{1+|\nabla u|^2}}\Bigr)_y\Bigr] \notag\\
 &=(1+u_y^2)u_{xx}+(1+u_x^2)u_{yy}-2u_xu_yu_{xy}.
\label{eq:1.7}
\end{align}

下面我们要说明: $\Omega$ 上满足极小曲面方程的函数, 它的图不仅是面积泛函的
临界点, 而且实际上在柱体 $\Omega\times\R\subset\R^3$ 中是面积最小的. 设
$\omega$ 是 $\Omega\times\R$ 上的二形式, 它由下式给出: 对 $X,Y\in\R^3$,
\begin{equation}
\omega(X,Y)=\det(X,Y,N),
\label{eq:1.8}
\end{equation}
其中
\begin{equation}
N=\frac{(-u_x,-u_y,1)}{\sqrt{1+|\nabla u|^2}}.
\label{eq:1.9}
\end{equation}
注意到
\begin{align}
\omega\Bigl(\frac{\partial}{\partial x},\frac{\partial}{\partial y}\Bigr)
 &=\frac{1}{\sqrt{1+|\nabla u|^2}}
   \begin{vmatrix} 1 & 0 & -u_x\\ 0 & 1 & -u_y\\ 0 & 0 & 1\end{vmatrix}
  =\frac{1}{\sqrt{1+|\nabla u|^2}},
\label{eq:1.10}\\[2pt]
\omega\Bigl(\frac{\partial}{\partial y},\frac{\partial}{\partial z}\Bigr)
 &=\frac{1}{\sqrt{1+|\nabla u|^2}}
   \begin{vmatrix} 0 & 0 & -u_x\\ 1 & 0 & -u_y\\ 0 & 1 & 1\end{vmatrix}
  =\frac{-u_x}{\sqrt{1+|\nabla u|^2}},
\label{eq:1.11}
\end{align}
以及
\begin{equation}
\omega\Bigl(\frac{\partial}{\partial x},\frac{\partial}{\partial z}\Bigr)
 =\frac{1}{\sqrt{1+|\nabla u|^2}}
  \begin{vmatrix} 1 & 0 & -u_x\\ 0 & 0 & -u_y\\ 0 & 1 & 1\end{vmatrix}
 =\frac{u_y}{\sqrt{1+|\nabla u|^2}}.
\label{eq:1.12}
\end{equation}
于是
\begin{equation}
\omega=\frac{dx\wedge dy-u_x\,dy\wedge dz-u_y\,dz\wedge dx}
            {\sqrt{1+|\nabla u|^2}},
\label{eq:1.13}
\end{equation}
并且
\begin{equation}
d\omega=\frac{\partial}{\partial x}
  \Bigl(\frac{-u_x}{\sqrt{1+|\nabla u|^2}}\Bigr)
 +\frac{\partial}{\partial y}
  \Bigl(\frac{-u_y}{\sqrt{1+|\nabla u|^2}}\Bigr)=0,
\label{eq:1.14}
\end{equation}
因为 $u$ 满足极小曲面方程. 总之, 形式 $\omega$ 是闭的, 并且对点 $(x,y,z)$
处的任意正交单位向量 $X$ 与 $Y$,
\begin{equation}
|\omega(X,Y)|\leq 1,
\label{eq:1.15}
\end{equation}
其中等号成立当且仅当
\begin{equation}
X,Y\subset T_{(x,y,u(x,y))}\Graph_u.
\label{eq:1.16}
\end{equation}
这样的形式 $\omega$ 称为\newterm{calibration}{校准}{calibration}, 可以用它
证明 $\Graph_u$ 是面积最小的:

\begin{lemma}
\label{lem:1.1}
若 $u:\Omega\to\R$ 满足极小曲面方程, 而 $\Sigma\subset\Omega\times\R$ 是
另一张满足 $\partial\Sigma=\partial\Graph_u$ 的曲面, 则
\begin{equation}
\Area(\Graph_u)\leq\Area(\Sigma).
\label{eq:1.17}
\end{equation}
\end{lemma}

\begin{proof}
由于 $\omega$ 是闭形式, 且 $\Graph_u$ 与 $\Sigma$ 同调, Stokes 定理给出
\begin{equation}
\int_{\Graph_u}\omega=\int_{\Sigma}\omega.
\label{eq:1.18}
\end{equation}
把它与 \eqref{eq:1.15}、\eqref{eq:1.16} 结合起来, 便得到
\begin{equation}
\Area(\Graph_u)=\int_{\Graph_u}\omega=\int_{\Sigma}\omega\leq\Area(\Sigma).
\label{eq:1.19}
\end{equation}
\end{proof}

本书中我们始终用 $B_r(x)$ 表示 $\R^3$(或 $\R^{n+1}$)中以 $x$ 为心、$r$ 为
半径的球; 常把 $B_r(0)$ 简记为 $B_r$. 用 $D_r$ 表示 $\R^2$ 中以 $0$ 为心、
$r$ 为半径的圆盘.

\begin{corollary}
\label{cor:1.2}
若 $u:\Omega\to\R$ 满足极小曲面方程, 且 $D_r\subset\Omega$, 则
\begin{equation}
\Area(B_r\cap\Graph_u)\leq
  \frac{\Area(\bS^2)}{2}\,r^2=2\pi r^2.
\label{eq:1.20}
\end{equation}
\end{corollary}

\begin{proof}
由于 $\partial B_r\cap\Graph_u$ 把 $\partial B_r$ 分成两部分, 其中至少有一
部分的面积不超过 $(\Area(\bS^2)/2)r^2$, 故由 \eqref{eq:1.17} 即得
\eqref{eq:1.20}.
\end{proof}

若区域 $\Omega$ 是凸的, 则极小图是绝对面积最小的. 为此先注意: 对凸集
$\Omega$, 最近点投影 $P:\R^3\to\Omega\times\R$ 是距离不增的 Lipschitz 映射,
且在 $\Omega\times\R$ 上是恒等映射. 若 $\Sigma\subset\R^3$ 是任一满足
$\partial\Sigma=\partial\Graph_u$ 的曲面, 则 $\Sigma'=P(\Sigma)$ 满足
$\Area(\Sigma')\leq\Area(\Sigma)$. 对 $\Sigma'$ 用 \eqref{eq:1.17}, 便知
$\Area(\Graph_u)\leq\Area(\Sigma')$, 结论成立.

与上面非常类似的计算表明: 若 $\Omega\subset\R^{n-1}$, $u:\Omega\to\R$ 是
$C^2$ 函数, 则 $u$ 的图是面积泛函的临界点, 当且仅当 $u$ 满足方程
\begin{equation}
\dive\Bigl(\frac{\nabla u}{\sqrt{1+|\nabla u|^2}}\Bigr)=0.
\label{eq:1.21}
\end{equation}
此外, 与 \eqref{eq:1.17} 一样, 若 $u$ 满足 \eqref{eq:1.21}, 则 $u$ 的图
实际上是面积最小的. 从而, 与 \eqref{eq:1.20} 一样, 若 $\Omega$ 含有半径为
$r$ 的球, 则
\begin{equation}
\Vol(B_r\cap\Graph_u)\leq
  \frac{\Vol(\bS^{n-1})}{2}\,r^{n-1}.
\label{eq:1.22}
\end{equation}

\mtpara{子流形的几何}
我们本来也可以更一般地考虑一个 $k$ 维子流形 $\Sigma$(可以带边), 它位于某个
Riemann 流形 $M$ 之中(度量为 $g$, 协变导数为 $\nabla$), 并且是面积泛函的
临界点.

以下若 $X$ 是 $\Sigma\subset M$ 上的向量场, 则用 $X^T$ 与 $X^N$ 分别表示其
切向分量与法向分量. $M$ 上的协变导数 $\nabla$ 诱导出 $\Sigma$ 上的协变导数
$\nabla_{\Sigma}$ 与 $\Sigma$ 的\newterm{second-fundamental-form}{第二基本
形式}{second fundamental form} $A$. 也就是说, 诱导协变导数由
\begin{equation}
\nabla_{\Sigma}=(\nabla)^T
\label{eq:1.23}
\end{equation}
给出, 而 $\Sigma$ 上取值于向量空间的二线性形式 $A$ 由下式给出: 对
$X,Y\in T_x\Sigma$,
\begin{equation}
A(X,Y)=(\nabla_XY)^N.
\label{eq:1.24}
\end{equation}
由于切向量场的括号仍是切向量场, 容易看出 $A$ 是对称的, 即 $A(X,Y)=A(Y,X)$.
注意到
\begin{align}
\sum_{\ell=1}^{n-k}g\bigl(A(X,Y),N_\ell\bigr)N_\ell
 &=\sum_{\ell=1}^{n-k}g\bigl(\nabla_XY,N_\ell\bigr)N_\ell \notag\\
 &=-\sum_{\ell=1}^{n-k}g\bigl(Y,\nabla_XN_\ell\bigr)N_\ell,
\label{eq:1.25}
\end{align}
其中 $N_\ell$ 是 $\Sigma$ 在 $x$ 的某邻域内法空间的单位正交向量场基.

点 $x$ 处的\newterm{mean-curvature}{平均曲率}{mean curvature}向量 $H$
按定义是
\begin{equation}
H=\sum_{i=1}^{k}A(E_i,E_i),
\label{eq:1.26}
\end{equation}
其中 $E_i$ 是 $T_x\Sigma$ 的单位正交基. 此外, 第二基本形式在 $x$ 处的模长
平方按定义是
\begin{equation}
|A|^2=\sum_{i,j=1}^{k}\bigl|A(E_i,E_j)\bigr|^2.
\label{eq:1.27}
\end{equation}
还要回忆 Gauss 方程: 若 $X,Y\in T_x\Sigma$, 则
\begin{align}
K_{\Sigma}(X,Y)|X\wedge Y|^2
 &=K_M(X,Y)|X\wedge Y|^2 \notag\\
 &\quad+g\bigl(A(X,X),A(Y,Y)\bigr)-g\bigl(A(X,Y),A(X,Y)\bigr),
\label{eq:1.28}
\end{align}
其中 $|X\wedge Y|^2$ 由
\begin{equation}
|X\wedge Y|^2=g(X,X)g(Y,Y)-g(X,Y)^2
\label{eq:1.29}
\end{equation}
给出, 而 $K_M(X,Y)$ 与 $K_{\Sigma}(X,Y)$ 分别是 $M$ 与 $\Sigma$ 在 $X$、$Y$
张成的二平面方向上的截面曲率. 若 $\Sigma^{n-1}\subset M^n$ 是超曲面, $N$ 是
$x$ 的某邻域内的单位法向量场, 则
\begin{equation}
\nabla_{(\cdot)}N:T_x\Sigma\to T_x\Sigma
\label{eq:1.30}
\end{equation}
是对称映射(常称为 Weingarten 映射), 它的特征值 $(\kappa_i)_{i=1,\dots,n-1}$
称为\newterm{principal-curvature}{主曲率}{principal curvatures}. 此外,
\begin{equation}
g(H,N)=-\sum_{i=1}^{n-1}\kappa_i.
\label{eq:1.31}
\end{equation}
最后, 若 $X$ 是定义在 $\Sigma$ 的某邻域内的向量场, 则 $X$ 在 $x\in\Sigma$
处的散度是
\begin{equation}
\dive_{\Sigma}X=\sum_{i=1}^{n-1}g\bigl(\nabla_{E_i}X,E_i\bigr),
\label{eq:1.32}
\end{equation}
其中 $E_i$ 是 $T_x\Sigma$ 的单位正交基. 注意 $\dive_{\Sigma}$ 满足 Leibniz
法则
\begin{equation}
\dive_{\Sigma}(fX)=\langle\nabla_{\Sigma}f,X\rangle
  +f\,\dive_{\Sigma}(X).
\label{eq:1.33}
\end{equation}
我们也可以用 $\dive_{\Sigma}$ 定义 $\Sigma$ 上的 Laplace 算子
$\Delta_{\Sigma}$:
\begin{equation}
\Delta_{\Sigma}f=\dive_{\Sigma}(\nabla_{\Sigma}f).
\label{eq:1.34}
\end{equation}
若 $\Delta_{\Sigma}f=0$, 则称函数 $f$ 在 $\Sigma$ 上是\newterm{harmonic}
{调和}{harmonic}的.

\begin{remark}
\label{rem:1.3}
注意
\begin{align}
\dive_{\Sigma}Y^N
 &=\sum_i g\bigl(E_i,\nabla_{E_i}Y^N\bigr)
  =-\sum_i g\bigl(Y^N,\nabla_{E_i}E_i\bigr) \notag\\
 &=-g\bigl(Y^N,H\bigr).
\label{eq:1.35}
\end{align}
\end{remark}

\mtpara{第一变分公式}
设 $F:\Sigma\times(-\varepsilon,\varepsilon)\to M$ 是 $\Sigma$ 的一个具有
紧支集且边界固定的变分. 也就是说, $F=\mathrm{Id}$ 在某个紧集之外成立,
\begin{equation}
F(x,0)=x,
\label{eq:1.36}
\end{equation}
并且对一切 $x\in\partial\Sigma$,
\begin{equation}
F(x,t)=x.
\label{eq:1.37}
\end{equation}
限制在 $\Sigma$ 上的向量场 $F_t$ 常称为\newterm{variational-vector-field}
{变分向量场}{variational vector field}. 现在我们要计算这一族单参数曲面的
面积的第一变分. 设 $x_i$ 是 $\Sigma$ 上的局部坐标, 令
\begin{equation}
g_{ij}(t)=g\bigl(F_{x_i},F_{x_j}\bigr),
\label{eq:1.38}
\end{equation}
\begin{equation}
\nu(t)=\sqrt{\det\bigl(g_{ij}(t)\bigr)}\sqrt{\det\bigl(g^{ij}(0)\bigr)},
\label{eq:1.39}
\end{equation}
其中 $a^{ij}$ 表示矩阵 $a_{ij}$ 的逆. 注意 $\nu(t)$ 是良定义的, 与 $\Sigma$
上坐标系的选择无关(因为 $\det(g_{ij}(t))$ 按坐标变换微分之行列式的平方
变化, 而 $\det(g^{ij}(0))$ 按它的逆变化). 此外, 面积公式是
\begin{equation}
\Vol\bigl(F(\Sigma,t)\bigr)=\int\nu(t)\sqrt{\det\bigl(g_{ij}(0)\bigr)},
\label{eq:1.40}
\end{equation}
其中积分在 $\Sigma$ 上进行. 对它求导得
\begin{equation}
\frac{d}{dt}\Big|_{t=0}\Vol\bigl(F(\Sigma,t)\bigr)
 =\int\frac{d}{dt}\Big|_{t=0}\nu(t)\sqrt{\det\bigl(g_{ij}(0)\bigr)}.
\label{eq:1.41}
\end{equation}
为计算某点 $x$ 处的 $\frac{d}{dt}|_{t=0}\nu(t)$, 我们可以选取坐标系使其在
$x$ 处是单位正交的, 即在点 $x$ 处
\begin{equation}
g_{ij}(0)=\delta_{ij}=
\begin{cases}
0, & \text{若 } i\neq j,\\
1, & \text{若 } i=j.
\end{cases}
\label{eq:1.42}
\end{equation}
利用这一点以及 $t$ 与 $x_i$ 的导数可交换(即
$\nabla_{F_t}F_{x_i}-\nabla_{F_{x_i}}F_t=[F_t,F_{x_i}]=0$), 我们在 $x$ 处得到
\begin{align}
\frac{d}{dt}\Big|_{t=0}\nu(t)
 &=\frac12\sum_{i=1}^{k}\frac{d}{dt}
   \langle F_{x_i},F_{x_i}\rangle
  =\sum_{i=1}^{k}\langle\nabla_{F_t}F_{x_i},F_{x_i}\rangle \notag\\
 &=\sum_{i=1}^{k}\langle\nabla_{F_{x_i}}F_t,F_{x_i}\rangle
  =\dive_{\Sigma}F_t.
\label{eq:1.43}
\end{align}
把向量场 $F_t$ 写成其法向部分与切向部分之和, 便可把这个公式与平均曲率联系
起来, 得到
\begin{align}
\frac{d}{dt}\Big|_{t=0}\nu(t)
 &=\sum_{\ell=1}^{n-k}\sum_{i=1}^{k}
   \bigl\langle\nabla_{F_{x_i}}\langle F_t,N_\ell\rangle N_\ell,F_{x_i}\bigr\rangle
   +\dive_{\Sigma}F_t^T \notag\\
 &=\sum_{\ell=1}^{n-k}\sum_{i=1}^{k}
   \langle F_t,N_\ell\rangle\langle\nabla_{F_{x_i}}N_\ell,F_{x_i}\rangle
   +\dive_{\Sigma}F_t^T \notag\\
 &=-\langle F_t,H\rangle+\dive_{\Sigma}F_t^T.
\label{eq:1.44}
\end{align}
这里 $N_\ell$ 是 $\Sigma$ 在 $x$ 处法丛的单位正交基. 对 \eqref{eq:1.43} 与
\eqref{eq:1.44} 积分, 便得到所谓的第一变分公式:
\begin{equation}
\frac{d}{dt}\Big|_{t=0}\Vol\bigl(F(\Sigma,t)\bigr)
 =-\int_{\Sigma}\langle F_t,H\rangle
 =\int_{\Sigma}\dive_{\Sigma}F_t.
\label{eq:1.45}
\end{equation}
注意这里用了 Stokes 定理来说明 $\int\dive_{\Sigma}F_t^T=0$. 由
\eqref{eq:1.45} 可知, $\Sigma$ 是面积泛函的临界点, 当且仅当平均曲率 $H$
恒等于零.

\begin{definition}
\label{def:1.4}
\emph{(极小子流形)} 若平均曲率 $H$ 恒等于零, 则称浸入子流形
$\Sigma^k\subset M^n$ 是\textbf{极小}的.
\end{definition}

由恒等式 \eqref{eq:1.45} 可知, $\R^3$ 中的图是极小曲面, 当且仅当它满足
极小曲面方程 \eqref{eq:1.6}.

% ================================================================
\section{$\R^3$ 中极小曲面的例子}
\label{sec:1.2}

极小曲面的特征是平均曲率为零. 最简单的例子是第二基本形式(而不只是它的迹)
处处为零的情形, 即曲面是\newterm{totally-geodesic}{全测地}{totally geodesic}
的. 由于第二基本形式是单位法向量的导数, 这意味着法向量是常向量, 从而曲面是
一张平面.

\mtpara{曲面的拓扑}
不带边的紧可定向曲面按\newterm{genus}{亏格}{genus}分类, 亏格是一个非负整数.
亏格 $=0$ 对应球面, 亏格 $=1$ 对应环面. 亏格 $=k$ 的曲面以球面添加 $k$ 个
环柄所成的曲面为模型. 带边的紧可定向曲面是把这样的曲面去掉若干互不相交的
圆盘而得到的. 带边曲面的亏格就是原曲面的亏格, 而边界对应于去掉圆盘后产生的
曲面边缘. 特别地, 亏格为 $0$ 且边界非空的曲面是\newterm{planar-domain}{平面区域}{planar domain}, 即它可以由平面上的圆盘去掉若干互不相交的子圆盘得到.
有时我们会谈到单连通的曲面, 意思是曲面上的每个闭路都可以(不离开曲面地)收缩
为一点. 容易看出, 单连通的曲面只有圆盘与球面.

亏格定义的一个直接推论是它在包含关系下的单调性:

\begin{lemma}
\label{lem:1.5}
若 $\Sigma$ 的亏格为 $k$, $\Sigma_0\subset\Sigma$, 则 $\Sigma_0$ 的亏格至
多为 $k$.
\end{lemma}

\mtpara{螺旋面; 见图~\ref{fig:1.1}}
\newterm{helicoid}{螺旋面}{helicoid}是集合 $x_3=\tan^{-1}(x_2/x_1)$;
等价地, 它由参数形式
\begin{equation}
(x_1,x_2,x_3)=(t\cos s,t\sin s,s)
\label{eq:1.46}
\end{equation}
给出, 其中 $s,t\in\R$. 它由 Meusnier(Monge 的学生)于 1776 年发现. 它是完全
的、嵌入的、单周期的、单连通的.

螺旋面是\newterm{ruled-surface}{直纹曲面}{ruled surface}, 因为它与水平面
$\{x_3=s\}$ 的交是直线. 这些直线以常速提升并旋转, 构成双螺旋阶梯. 1842 年
Catalan 证明螺旋面是唯一的(非平坦)直纹极小曲面; 关于进一步的唯一性结果见
第 8 章. 若曲面可参数化为
\begin{equation}
X(s,t)=\beta(t)+s\,\delta(t)\quad\text{其中 } s,t\in\R,
\label{eq:1.47}
\end{equation}
且 $\beta$ 与 $\delta$ 是 $\R^3$ 中的曲线, 则称该曲面是``直纹''的. 曲线
$\beta(t)$ 称为曲面的\newterm{directrix}{准线}{directrix}, 以 $\delta(t)$ 为
方向向量的直线称为一条\newterm{ruling}{直母线}{ruling}. 对标准螺旋面,
$x_3$ 轴是一条准线, 而对每个固定的 $t$, 直线 $s\mapsto(s\cos t,s\sin t,t)$
是一条直母线.

\mtfigure[0.32]{fig-1-1}{螺旋面, 图中画出了直母线. 图片提供: Matthias Weber,
  \url{www.indiana.edu/~minimal}.}{fig:1.1}

\mtpara{悬链面}
\newterm{catenoid}{悬链面}{catenoid}(图~\ref{fig:1.2})是唯一的非平坦旋转
极小曲面. 它由 Euler 于 1744 年发现, 并由 Meusnier(Monge 的学生)于 1776 年
证明是极小的. 它是完全的嵌入拓扑圆环(即亏格为零、两个端), 由 $\R^3$ 中
$x_1^2+x_2^2=\cosh^2(x_3)$ 的集合给出. 容易看出悬链面具有有限总曲率.

\mtfigure[0.55]{fig-1-2}{悬链面: 把 $x_1=\cosh x_3$ 绕 $x_3$ 轴旋转所得.
  图片提供: Matthias Weber, \url{www.indiana.edu/~minimal}.}{fig:1.2}

\mtpara{Scherk 的双周期曲面}
大约 70 年之后, 尽管 Monge、Legendre 与 Poisson 做了大量工作\footnote{关于
这段历史, 见 \cite{Os5}.}, Scherk \cite{Sk} 于 1835 年发现了下一个完全极小
曲面.

Scherk 曲面是双周期极小曲面, 它被定义为棋盘格中白色方格上方的图, 方格角点
处立有竖直线. Scherk 曲面的一个基本域是正方形 $|x_1|<\pi/2$、
$|x_2|<\pi/2$ 上的图
\begin{equation}
x_3=\log\frac{\cos(x_2)}{\cos(x_1)}.
\label{eq:1.48}
\end{equation}
若按平移取商, 得到 $T^2\times\R$ 中的极小曲面, 所得曲面的亏格为 $0$, 且有
$4$ 个环形端.

我们来验证 Scherk 曲面确实是极小曲面. 只需验证
$x_3=x_3(x_1,x_2)$ 满足极小曲面方程. 显然
\begin{align*}
\partial_{x_1}x_3&=\tan(x_1), &
\partial_{x_2}x_3&=-\tan(x_2),\\
\partial_{x_1x_1}x_3&=1+\tan^2(x_1), &
\partial_{x_2x_2}x_3&=-1-\tan^2(x_2),\\
\partial_{x_1x_2}x_3&=0. &&
\end{align*}
因而 $x_3$ 满足极小曲面方程 \eqref{eq:1.7}.

\mtfigure[0.45]{fig-1-3}{Scherk 的双周期曲面. 图片提供: Matthias Weber,
  \url{www.indiana.edu/~minimal}.}{fig:1.3}

\mtpara{Scherk 的单周期曲面}
同样在 1835 年, Scherk \cite{Sk} 发现了一个单周期极小曲面, 它看起来像一对
正交平面的去奇点化. 这个曲面至今仍是最重要的例子之一, 因为它在极小曲面与
常平均曲率曲面的粘合与去奇点构造中起关键作用; 见 Kapouleas \cite{Ka} 与
Traizet \cite{Tr1}.

\mtfigure[0.32]{fig-1-4}{Scherk 的单周期曲面. 图片提供: Matthias Weber,
  \url{www.indiana.edu/~minimal}.}{fig:1.4}

\mtpara{Riemann 例子}
约 1860 年, Riemann \cite{Ri} 分类了 $\R^3$ 中所有被水平面内的圆与直线
叶状化的极小曲面. 他证明这样的曲面只有平面、悬链面、螺旋面, 以及一个双参数
族, 后者现在称为 \newterm{riemann-examples}{Riemann 例子}{Riemann examples}.
他发现的是一族完全的嵌入极小曲面, 它们是单周期的且亏格为零. 每个曲面都有
无穷多个平行的平面端, 各由颈部(``裤子形曲面'')相连.

在刚体运动的意义下, 这是一个双参数极小曲面族. 参数为:
\begin{itemize}
  \item 颈部的大小;
  \item 周期向量与端之间的夹角.
\end{itemize}

\mtfigure[0.45]{fig-1-5}{Riemann 例子. 图片提供: Matthias Weber,
  \url{www.indiana.edu/~minimal}.}{fig:1.5}

若固定颈部大小而让夹角趋于竖直(即垂直于平面端), 该族退化为两个定向相反的
螺旋面. 另一方面, 当夹角趋于零时, 该族退化为悬链面.

\mtfigure[0.30]{fig-1-6}{Riemann 例子: 开始退化为螺旋面. 图片提供:
  Matthias Weber, \url{www.indiana.edu/~minimal}.}{fig:1.6}

\mtpara{Enneper 曲面}
1864 年, Enneper \cite{E} 发现了一个由
\begin{equation}
(x_1,x_2,x_3)=\Bigl(s-\frac{s^3}{3}+st^2,\;-t-s^2t+\frac{t^3}{3},\;s^2-t^2\Bigr)
\label{eq:1.49}
\end{equation}
参数化的极小曲面, 其中 $s,t\in\R$. 与前面的例子不同, Enneper 曲面不是嵌入
的. 与悬链面一样, Enneper 曲面具有有限总曲率.

\mtfigure[0.42]{fig-1-7}{Enneper 曲面中一小块嵌入的部分; 延展开去它会自相交.
  图片提供: Matthias Weber, \url{www.indiana.edu/~minimal}.}{fig:1.7}

\mtpara{Costa--Hoffman--Meeks 曲面}
1984 年, C. Costa \cite{Cc} 发现了一个新的极小曲面, 它具有有限总曲率、亏格
一、三个端, 而且每个端都是嵌入的. 这个曲面是用 Weierstrass 表示、通过在环面
上巧妙地选取数据构造出来的. 1985 年, D. Hoffman 与 W. Meeks \cite{HoMe1}
证明 Costa 曲面是嵌入的; 这个曲面现在称为
\newterm{chm-surface}{Costa--Hoffman--Meeks 曲面}{Costa--Hoffman--Meeks
surface}.

此外, Hoffman 与 Meeks 还证明 Costa 曲面只是一个具有三个端、且亏格递增的
曲面族中的第一个. 当亏格趋于无穷时, 这些曲面(至少作为集合)收敛到一张平面与
一个悬链面之并.

\mtfigure[0.40]{fig-1-8}{原始的 Costa--Hoffman--Meeks 曲面. 图片提供:
  Matthias Weber, \url{www.indiana.edu/~minimal}.}{fig:1.8}

\mtfigure[0.45]{fig-1-9}{有许多环柄的 Costa--Hoffman--Meeks 曲面. 图片提供:
  Matthias Weber, \url{www.indiana.edu/~minimal}.}{fig:1.9}

\mtfigure[0.48]{fig-1-10}{当环柄数很大时, 可以看到 Scherk 单周期曲面的弯曲
  版本. 图片提供: Matthias Weber, \url{www.indiana.edu/~minimal}.}{fig:1.10}

\mtpara{Weber--Wolf 曲面}
M. Weber 与 M. Wolf \cite{WeWo2} 构造了一族嵌入极小曲面: 它们有两个悬链面
型的端(上端与下端), 中间有任意多个平面端, 而环柄数至少与平面端数一样多.
人们猜测这些曲面正处在可能情形的边界上:

\begin{conjecture}[Hoffman--Meeks]
\label{conj:1.6}
若 $\Sigma$ 是 $\R^3$ 中具有 $k$ 个端、亏格为 $g$ 的完全嵌入极小曲面, 则
$k\leq g+2$.
\end{conjecture}

\mtfigure[0.42]{fig-1-11}{亏格为 4、有 5 个端的 Weber--Wolf 曲面. 图片提供:
  Matthias Weber, \url{www.indiana.edu/~minimal}.}{fig:1.11}

\mtpara{亏格一螺旋面}
1993 年, Hoffman--Karcher--Wei 给出了一个渐近于螺旋面的完全嵌入亏格一极小
曲面存在的数值证据, 并称之为``亏格一螺旋面''. 在 \cite{HoWW} 中, Hoffman、
Weber 与 Wolf 把这样的曲面构造为``单周期亏格一螺旋面''的极限, 其中每个
单周期亏格一螺旋面都是通过 Weierstrass 表示构造的. 后来, Hoffman 与 White
在 \cite{HoWh1} 中用变分方法构造了亏格一螺旋面.

\mtfigure[0.42]{fig-1-12}{亏格一螺旋面. 图片提供: Matthias Weber,
  \url{www.indiana.edu/~minimal}.}{fig:1.12}

\mtfigure[0.28]{fig-1-13}{一个渐近于螺旋面的周期极小曲面, 其基本域的亏格为
  一. 图片提供: Matthias Weber, \url{www.indiana.edu/~minimal}.}{fig:1.13}

\mtpara{嵌入曲面的构件}
Riemann 例子、亏格一螺旋面与 Costa--Hoffman--Meeks 曲面都是亏格有限的嵌入
极小曲面, 每一个都可以看作由平面、螺旋面与悬链面的片段拼成的. 在第 8 章中,
我们将讨论关于嵌入极小曲面的 Colding--Minicozzi 关键结构结果. 这个结果来自
\cite{CM3}--\cite{CM7}, 它证明了: 任何亏格有限的嵌入极小曲面都可以由平面、
螺旋面与悬链面拼成.

\mtpara{Schwarz 的三周期极小曲面}
19 世纪后期, H. A. Schwarz(第 5 章讨论的 Schwarz 交替法即出自他)及其学生,
包括 Neovius \cite{Ne}, 在 $\R^3$ 中发现了五个三周期极小曲面; 见 \cite{Sz2},
\cite{Sz3}. 这些曲面的亏格都是无穷, 但它们关于晶体群的商的亏格有限. 1970 年,
A. Schoen \cite{Sca} 构造了许多新的三周期极小曲面族, 其中包括陀螺面
(\newterm{gyroid}{陀螺面}{gyroid}). 这些曲面的存在性由 Karcher 于 1989 年
\cite{Kr} 严格证明. 陀螺面的嵌入性由 K. Grosse-Brauckmann 与 M. Wohlgemuth
于 1996 年 \cite{GrW} 证明. 三周期极小曲面在科学中有应用; 见 \cite{AHHT}.
例如, 陀螺面结构出现在某些嵌段共聚物中.

\mtfigure[0.45]{fig-1-14}{Schwarz D 曲面. 图片提供: Matthias Weber,
  \url{www.indiana.edu/~minimal}.}{fig:1.14}

\mtfigure[0.45]{fig-1-15}{Schwarz P 曲面. 图片提供: Matthias Weber,
  \url{www.indiana.edu/~minimal}.}{fig:1.15}

Schwarz D 曲面的基本域发现于 1865 年, 它解决了一个 Plateau 问题, 其边界位于
一个立方体的棱上. Schwarz P 曲面的基本域发现于 1880 年代, 它同样可以通过解
Plateau 问题来构造, 这次的边界是一个四边形, 其顶点是正八面体的顶点. 这种
``用多边形边界解 Plateau 问题再作反射''的想法已被非常成功地运用, 包括
H. B. Lawson 在圆三维球中构造极小曲面 \cite{La1}; 参看 6.2 小节.

\mtfigure[0.62]{fig-1-16}{Neovius 极小曲面的一个基本域. 图片提供:
  Matthias Weber, \url{www.indiana.edu/~minimal}.}{fig:1.16}

\mtfigure[0.45]{fig-1-17}{三周期的 Neovius 曲面, 发现于 1883 年. 图片提供:
  Matthias Weber, \url{www.indiana.edu/~minimal}.}{fig:1.17}

\mtfigure[0.48]{fig-1-18}{Schwarz CLP 曲面的一个基本域. 图片提供:
  Matthias Weber, \url{www.indiana.edu/~minimal}.}{fig:1.18}

\mtfigure[0.45]{fig-1-19}{三周期的 CLP 曲面. 图片提供: Matthias Weber,
  \url{www.indiana.edu/~minimal}.}{fig:1.19}

\mtfigure[0.55]{fig-1-20}{Schwarz H 曲面的一个基本域. 图片提供:
  Matthias Weber, \url{www.indiana.edu/~minimal}.}{fig:1.20}

\mtfigure[0.45]{fig-1-21}{三周期的 Schwarz H 曲面. 图片提供: Matthias Weber,
  \url{www.indiana.edu/~minimal}.}{fig:1.21}

\mtpara{Chen--Gackstatter 曲面}
1981 年, Chen 与 Gackstatter \cite{CcG} 发现了第一个亏格为 1、具有有限总
曲率的完全浸入极小曲面. Chen--Gackstatter 曲面有一个端, 且是 Enneper 型的.

M. Weber 与 M. Wolf \cite{WeWo1} 证明, 可以给 Chen--Gackstatter 曲面添加
环柄, 得到亏格任意高、带一个 Enneper 端的曲面. S. Kim 在其学位论文 \cite{Ks}
中证明: 当环柄数趋于无穷时, Weber--Wolf 例子收敛到 Scherk 单周期曲面.

\mtfigure[0.45]{fig-1-22}{Chen--Gackstatter 曲面. 图片提供: Matthias Weber,
  \url{www.indiana.edu/~minimal}.}{fig:1.22}

\mtfigure[0.45]{fig-1-23}{给 Chen--Gackstatter 曲面添加环柄; 它越来越像
  Scherk 单周期曲面. 图片提供: Matthias Weber,
  \url{www.indiana.edu/~minimal}.}{fig:1.23}

\mtpara{高维悬链面}
悬链面有一个高维类比, 即绕一条轴旋转的极小旋转曲面. 关于它的主要性质我们
引用 \cite{TZ} 的第 2 节. 从我们的角度看, 最重要的几条是:
\begin{itemize}
  \item 它是完全的嵌入极小超曲面.
  \item $n$ 维悬链面由一条曲线绕轴旋转给出, 因而拓扑上是
    $\R\times\bS^{n-1}$. 特别地, 当 $n>2$ 时它是单连通的.
  \item $n$ 维悬链面在半径为 $R$ 的球中的体积不超过 $B_R\subset\R^n$ 体积
    的两倍, 并且当 $R$ 趋于无穷时收敛到该值.
  \item 当 $n>3$ 时, 悬链面位于两个平行超平面之间的一个板层中.
\end{itemize}

% ================================================================
\section{第一变分公式的推论}
\label{sec:1.3}

本节中, 我们将收集第一变分公式的一些重要推论. 其中最重要的是单调性公式, 即
命题~\ref{prop:1.12}. 在后面的章节中, 我们将回到这一主题: 在第 3 章把这些
结果推广到驻定\newterm{varifold}{变分流形}{varifold}, 在第 5 章推广到三维
流形中的极小曲面.

由 (1.45) 可以看出, $\Sigma$ 是极小的, 当且仅当对一切具有紧支集且在
$\Sigma$ 的边界上为零的向量场 $X$, 都有
\begin{equation}
\int_{\Sigma}\dive_{\Sigma}X=0.
\label{eq:1.50}
\end{equation}
这个方程称为\newterm{first-variation}{第一变分公式}{first variation formula}.
它的好处在于: 只要能在 $\Sigma$ 上定义散度, (1.50) 就有意义. (这使我们以后
能为极小曲面定义``弱解''的概念.)

本节余下部分中, 若 $x_0\in\R^n$ 固定, 则记 $B_s=B_s(x_0)$ 为以 $x_0$ 为心、
半径为 $s$ 的欧氏球.

\mtpara{坐标函数的调和性}
作为 (1.50) 的推论, 我们将证明下述命题:

\begin{proposition}
\label{prop:1.7}
$\Sigma^k\subset\R^n$ 是极小的, 当且仅当 $\R^n$ 的坐标函数限制在 $\Sigma$
上都是调和函数.
\end{proposition}

\begin{proof}
设 $e_i=\nabla_{\R^n}x_i$ 是 $\R^n$ 上的第 $i$ 个坐标向量场, $\eta$ 是
$\Sigma$ 上具有紧支集的光滑函数, 且在 $\partial\Sigma$ 上为零(若 $\Sigma$
有边界). 由于 $e_i$ 是常向量场, 紧支集向量场 $\eta e_i$ 的散度为
\begin{equation}
\dive_{\Sigma}(\eta e_i)=\langle\nabla_{\Sigma}\eta,e_i\rangle
 =\langle\nabla_{\Sigma}\eta,\nabla_{\Sigma}x_i\rangle,
\label{eq:1.51}
\end{equation}
于是
\begin{equation}
-\int_{\Sigma}\langle\eta e_i,H\rangle
 =\int_{\Sigma}\dive_{\Sigma}(\eta e_i)
 =\int_{\Sigma}\langle\nabla_{\Sigma}\eta,\nabla_{\Sigma}x_i\rangle
 =-\int_{\Sigma}\eta\,\Delta_{\Sigma}x_i.
\label{eq:1.52}
\end{equation}
由此即得结论.
\end{proof}

作为推论, 我们可以从 $(k-1)$ 维同调类出发, 定义一个到 $\R^n$ 的通量同态:

\begin{corollary}
\label{cor:1.8}
若 $\Sigma^k\subset\R^n$ 极小, 则存在定义良好的``通量同态'' $F$, 它把
$(k-1)$ 维同调类映到 $\R^n$, 由
\begin{equation}
F_i([\gamma])=\int_{\gamma}\langle n_{\gamma},e_i\rangle
\label{eq:1.53}
\end{equation}
给出, 其中 $n_{\gamma}$ 是 $\gamma$ 的定向余法向量\footnote{余法向量是指与
$\Sigma$ 相切、但与 $\gamma$ 正交的向量.}.
\end{corollary}

\begin{proof}
该映射显然是线性的, 故只需证明它只依赖于 $\gamma$ 的同调类. 然而, 若
$\gamma_1$ 与 $\gamma_2$ 同调, 则它们合起来在 $\Sigma$ 中张成一个 $k$-链
$\Omega$. 于是极小性与散度定理给出
\begin{equation}
\int_{\gamma_1}\langle n_{\gamma_1},e_i\rangle
 -\int_{\gamma_2}\langle n_{\gamma_2},e_i\rangle
 =\int_{\Omega}\dive_{\Sigma}e_i=0,
\label{eq:1.54}
\end{equation}
证毕.
\end{proof}

(1.50) 的另一个非常有用的推论, 是 $k$ 维极小曲面 $\Sigma$ 上 $|x|^2$ 的
Laplace 的公式. 即: 由于对任意向量 $v$ 都有
\begin{equation}
\nabla_v x_i=\langle v,e_i\rangle,
\label{eq:1.55}
\end{equation}
故
\begin{equation}
\dive_{\Sigma}(x_1,\dots,x_n)=k.
\label{eq:1.56}
\end{equation}
把 (1.56) 与 (1.50) 结合起来(其实我们只用到对任意 $Y$ 都有
$\dive_{\Sigma}Y^{N}=0$), 便得到
\begin{equation}
\Delta_{\Sigma}|x|^2
 =2\dive_{\Sigma}(x_1,\dots,x_n)^{T}
 =2\dive_{\Sigma}(x_1,\dots,x_n)=2k.
\label{eq:1.57}
\end{equation}

回忆: 若 $\Xi\subset\R^n$ 是紧子集, 则包含 $\Xi$ 的最小凸集(凸包
$\operatorname{Conv}(\Xi)$)是包含 $\Xi$ 的所有半空间之交. 在 \cite{Os6} 中,
Osserman 证明了极大值原理迫使极小子流形落在其边界的凸包之中:

\begin{proposition}[凸包性质, \cite{Os6}]
\label{prop:1.9}
若 $\Sigma^k\subset\R^n$ 是紧极小子流形, 则
$\Sigma\subset\operatorname{Conv}(\partial\Sigma)$.
\end{proposition}

\begin{proof}
$\R^n$ 中的半空间 $H$ 可写成
\begin{equation}
H=\{x\in\R^n\mid\langle x,e\rangle\leq a\},
\label{eq:1.58}
\end{equation}
其中 $e\in\bS^{n-1}$, $a\in\R$ 为常数. 由命题~\ref{prop:1.7}, 函数
$u(x)=\langle e,x\rangle$ 在 $\Sigma$ 上调和, 从而由极大值原理, 它在
$\partial\Sigma$ 上取到最大值.
\end{proof}

凸包性质证明中的论证可以换一种说法: 当我们把一个超平面朝极小曲面平移时,
第一个接触点必定落在边界上.

凸包性质给出极小子流形在拓扑上的单调性. 最简单的版本是关于极小圆盘的:

\begin{corollary}
\label{cor:1.10}
若 $\Sigma^2\subset\R^n$ 是紧极小圆盘, $K$ 是紧凸集且
$K\cap\partial\Sigma=\emptyset$, 则 $K\cap\Sigma$ 单连通.
\end{corollary}

\begin{proof}
否则, 存在闭曲线 $\gamma\subset K\cap\Sigma$, 它在 $K\cap\Sigma$ 中不围成
圆盘. 由于 $\Sigma$ 单连通, $\gamma$ 在 $\Sigma$ 中围成一个圆盘
$\Gamma\subset\Sigma$. 但此时极小圆盘 $\Gamma$ 违反凸包性质: 它的边界
$\partial\Gamma=\gamma$ 落在 $K$ 内, 而 $\Gamma$ 不落在 $K$ 内.
\end{proof}

更一般地, 可以得到 $(k-1)$ 阶同调群的单调性:

\begin{lemma}
\label{lem:1.11}
若 $\Sigma^k\subset\R^n$ 是紧极小子流形, $K$ 是紧凸集且
$K\cap\partial\Sigma=\emptyset$, 则 $K\cap\Sigma$ 到 $\Sigma$ 的包含映射在
$(k-1)$ 阶同调群上是单射.
\end{lemma}

\begin{proof}
否则, $B_R(0)\cap\Sigma$ 含有一个 $(k-1)$-闭链 $\gamma$, 它在
$B_R(0)\cap\Sigma$ 中不围成链, 但在 $\Sigma$ 中围成 $\Gamma\subset\Sigma$.
然而 $\Gamma$ 与凸包性质矛盾.
\end{proof}

\mtpara{单调性}
在陈述并证明极小子流形的体积单调性公式之前, 我们需要回忆余面积公式. 该公式
断言(证明可参看例如 \cite{Fe} 或 \cite{KP}): 若 $\Sigma$ 是流形,
\[
h\colon\Sigma\longrightarrow\R
\]
是 $\Sigma$ 上的固有(即对一切 $t\in\R$, $h^{-1}((-\infty,t])$ 紧)
Lipschitz 函数, 则对 $\Sigma$ 上一切局部可积函数 $f$ 及 $t\in\R$,
\begin{equation}
\int_{\{h\leq t\}}f\,|\nabla_{\Sigma}h|
 =\int_{-\infty}^{t}\int_{h=\tau}f\,d\tau.
\label{eq:1.59}
\end{equation}

\begin{proposition}[单调性公式]
\label{prop:1.12}
设 $\Sigma^k\subset\R^n$ 是极小子流形, $x_0\in\R^n$; 则对一切 $0<s<t$,
\begin{equation}
\begin{aligned}
&\frac{\Vol(B_t(x_0)\cap\Sigma)}{t^k}
 -\frac{\Vol(B_s(x_0)\cap\Sigma)}{s^k}\\
&\qquad=\int_{(B_t(x_0)\setminus B_s(x_0))\cap\Sigma}
 \frac{\bigl|(x-x_0)^{N}\bigr|^2}{|x-x_0|^{k+2}}.
\end{aligned}
\label{eq:1.60}
\end{equation}
\end{proposition}

\begin{proof}
在 $\Sigma$ 上定义函数 $f(x)=|x-x_0|$. 由于 $\Sigma$ 极小,
\begin{equation}
\Delta_{\Sigma}f^2=2\dive_{\Sigma}(x-x_0)=2k.
\label{eq:1.61}
\end{equation}
由 Stokes 定理, 对此积分得
\begin{equation}
2k\Vol(\{f\leq s\})=\int_{\{f\leq s\}}\Delta_{\Sigma}f^2
 =2\int_{\{f=s\}}\bigl|(x-x_0)^{T}\bigr|.
\label{eq:1.62}
\end{equation}
余面积公式(即 \eqref{eq:1.59})给出
\begin{equation}
\Vol(\{f\leq s\})
 =\int_{\{f\leq s\}}|\nabla_{\Sigma}f|^{-1}|\nabla_{\Sigma}f|
 =\int_0^s\int_{f=\tau}|\nabla_{\Sigma}f|^{-1}\,d\tau.
\label{eq:1.63}
\end{equation}
把它与 \eqref{eq:1.62} 结合, 经简单计算便得
\begin{equation}
\begin{aligned}
\frac{d}{ds}\bigl(s^{-k}\Vol(\{f\leq s\})\bigr)
&=-ks^{-k-1}\Vol(\{f\leq s\})
 +s^{-k}\int_{\{f=s\}}\frac{|x-x_0|}{|(x-x_0)^{T}|}\\
&=s^{-k-1}\int_{\{f=s\}}
 \Bigl(\frac{|x-x_0|^2}{|(x-x_0)^{T}|}-|(x-x_0)^{T}|\Bigr)\\
&=s^{-k-1}\int_{\{f=s\}}
 \frac{\bigl|(x-x_0)^{N}\bigr|^2}{|(x-x_0)^{T}|}.
\end{aligned}
\label{eq:1.64}
\end{equation}
再积分一次并应用余面积公式, 即得结论.
\end{proof}

注意, $(x-x_0)^{N}$ 恰好当 $\Sigma$ 关于 $x_0$ 为锥面时为零, 即当 $\Sigma$
在以 $x_0$ 为中心的伸缩下不变时为零. 作为推论, 我们得到:

\begin{corollary}
\label{cor:1.13}
设 $\Sigma^k\subset\R^n$ 是极小子流形, $x_0\in\R^n$; 则函数
\begin{equation}
\Theta_{x_0}(s)
 =\frac{\Vol(B_s(x_0)\cap\Sigma)}{\Vol(B_s\subset\R^k)}
\label{eq:1.65}
\end{equation}
是 $s$ 的非降函数. 此外, $\Theta_{x_0}(s)$ 关于 $s$ 为常数, 当且仅当
$\Sigma$ 关于 $x_0$ 为锥面.

最后, 若 $x_0\in\Sigma$, 则 $\Theta_{x_0}(s)\geq1$; 若对某个 $s>0$ 有
$\Theta_{x_0}(s)=1$, 则 $B_s\cap\Sigma$ 是某个 $k$ 维平面中的球.
\end{corollary}

\begin{proof}
命题~\ref{prop:1.12} 直接表明 $\Theta_{x_0}(s)$ 单调非降. 由于 $\Sigma$
光滑且固有, 它在无穷小意义下是欧氏的, 故
\begin{equation}
\lim_{s\to0}\Theta_{x_0}(s)\geq1.
\label{eq:1.66}
\end{equation}
把 $\Theta_{x_0}(s)$ 的单调性与 \eqref{eq:1.66} 结合, 得 $\Theta_{x_0}(s)\geq1$.
若有 $\Theta_{x_0}(s)=1$, 则 $\Theta_{x_0}$ 关于 $s$ 为常数, 从而由
\eqref{eq:1.60}, $(x-x_0)^{N}$ 恒为零. 显然这意味着 $\Sigma$ 在伸缩下不变; 又
$\Sigma$ 光滑, 故 $\Sigma$ 含于某个 $k$ 维平面之中.
\end{proof}

为后文引用方便, 我们记下推论~\ref{cor:1.13} 的几个推论. 设 $\Sigma$ 是极小
子流形, 定义 $x_0$ 处的密度为
\begin{equation}
\Theta_{x_0}=\lim_{s\to0}\Theta_{x_0}(s).
\label{eq:1.67}
\end{equation}
由于 $\Theta_{x_0}(s)$ 单调, 该极限存在; 且由 \eqref{eq:1.66}, 当 $x_0\in\Sigma$ 时
它总不小于 $1$. 事实上, 只要 $\Sigma$ 光滑, $\Theta_{x_0}$ 就是非负整数,
等于 $\Sigma$ 在 $x_0$ 处的重数. 注意若 $\Sigma$ 不是嵌入的, 则这个重数可以
大于 $1$.

下一个结果是单调性的一个初等推论, 它表明这个重数是上半连续的.

\begin{corollary}
\label{cor:1.14}
若 $\Sigma^k\subset\R^n$ 是极小子流形, 则密度 $\Theta_x$ 是 $\R^n$ 上的
上半连续函数. 因此, 对任意 $\Lambda\geq0$, 集合
\begin{equation}
\{x\in\Sigma\mid\Theta_x\geq\Lambda\}
\label{eq:1.68}
\end{equation}
是闭集.
\end{corollary}

\begin{proof}
我们要证明: 若 $x_j$ 是趋于 $x$ 的点列, 则
\begin{equation}
\Theta_x\geq\limsup_{x_j\to x}\Theta_{x_j}.
\label{eq:1.69}
\end{equation}
给定任意 $\delta>0$, 存在 $s>0$ 使得
\begin{equation}
\Theta_x\geq\Theta_x(2s)-\delta,
\label{eq:1.70}
\end{equation}
并且可取 $0<\epsilon<s$ 使得
\begin{equation}
\Theta_x\geq(1+s^{-1}\epsilon)^k\Theta_x(2s)-2\delta.
\label{eq:1.71}
\end{equation}
对任意满足 $|x-x_j|<\epsilon$ 的 $x_j$,
\begin{equation}
\begin{aligned}
\Theta_{x_j}&\leq\Theta_{x_j}(s)
 \leq\frac{\Vol(B_{s+\epsilon}(x)\cap\Sigma)}{\Vol(B_s\subset\R^k)}
 =(1+s^{-1}\epsilon)^k\Theta_x(s+\epsilon)\\
&\leq2\delta+\Theta_x,
\end{aligned}
\label{eq:1.72}
\end{equation}
其中最后一个不等式由 \eqref{eq:1.71} 得到. 由于 $\delta$ 可任意小, \eqref{eq:1.72} 蕴含
\eqref{eq:1.69}, 故 $\Theta$ 上半连续. 由此立即推出 \eqref{eq:1.68} 中定义的集合必为闭集.
\end{proof}

\mtpara{平均值不等式}

\begin{proposition}[平均值不等式]
\label{prop:1.15}
若 $\Sigma^k\subset\R^n$ 是极小子流形, $f$ 是 $\Sigma$ 上的函数, 则
\begin{equation}
\begin{aligned}
&t^{-k}\int_{B_t\cap\Sigma}f-s^{-k}\int_{B_s\cap\Sigma}f\\
&\quad=\int_{(B_t\setminus B_s)\cap\Sigma}f\,
 \frac{|x^{N}|^{2}}{|x|^{k+2}}
 +\frac12\int_s^t\tau^{-k-1}\int_{B_\tau\cap\Sigma}
 (\tau^2-|x|^2)\,\Delta_{\Sigma}f\,d\tau.
\end{aligned}
\label{eq:1.73}
\end{equation}
\end{proposition}

\begin{proof}
注意单调性公式正是 $f=1$ 的特殊情形. 由于 $\Sigma$ 极小, 分部积分给出
\begin{equation}
\begin{aligned}
2k\int_{B_s\cap\Sigma}f
&=\int_{B_s\cap\Sigma}f\,\Delta_{\Sigma}|x|^2\\
&=\int_{B_s\cap\Sigma}|x|^2\,\Delta_{\Sigma}f
 +2\int_{\partial B_s\cap\Sigma}f\,|x^{T}|
 -s^2\int_{B_s\cap\Sigma}\Delta_{\Sigma}f.
\end{aligned}
\label{eq:1.74}
\end{equation}
利用此式与余面积公式(即 \eqref{eq:1.59}), 得
\begin{equation}
\begin{aligned}
&\frac{d}{ds}\Bigl(s^{-k}\int_{B_s\cap\Sigma}f\Bigr)\\
&\quad=-ks^{-k-1}\int_{B_s\cap\Sigma}f
 +s^{-k}\int_{\partial B_s\cap\Sigma}f\,\frac{|x|}{|x^{T}|}\\
&\quad=s^{-k-1}\int_{\partial B_s\cap\Sigma}f\,
 \frac{|x^{N}|^{2}}{|x^{T}|}
 +\frac12s^{-k-1}\int_{B_s\cap\Sigma}
 (s^2-|x|^2)\,\Delta_{\Sigma}f.
\end{aligned}
\label{eq:1.75}
\end{equation}
积分并利用余面积公式, 即得结论.
\end{proof}

为以后引用, 我们接着记下由命题~\ref{prop:1.15} 得到的一个一般平均值不等式.

\begin{corollary}
\label{cor:1.16}
设 $\Sigma^k\subset\R^n$ 是极小子流形, $x_0\in\Sigma$, $s>0$ 满足
$B_s(x_0)\cap\partial\Sigma=\emptyset$. 若 $f$ 是 $\Sigma$ 上的非负函数, 且
$\Delta_{\Sigma}f\geq-\lambda s^{-2}f$, 则
\begin{equation}
f(x_0)\leq\mathrm{e}^{\lambda/2}
 \frac{\int_{B_s(x_0)\cap\Sigma}f}{\Vol(B_s\subset\R^k)}.
\label{eq:1.76}
\end{equation}
\end{corollary}

\begin{proof}
定义
\begin{equation}
g(t)=t^{-k}\int_{B_t(x_0)\cap\Sigma}f,
\label{eq:1.77}
\end{equation}
则命题~\ref{prop:1.15} 蕴含
\begin{equation}
g'(t)\geq-\frac{\lambda}{2}s^{-2}t^{1-k}
 \int_{B_t(x_0)\cap\Sigma}f
 =-\frac{\lambda}{2}s^{-2}t\,g(t).
\label{eq:1.78}
\end{equation}
可将 \eqref{eq:1.78} 改写为
\begin{equation}
\frac{g'(t)}{g(t)}\geq-\frac{\lambda}{2}s^{-2}t
 \geq-\frac{\lambda}{2s}.
\label{eq:1.79}
\end{equation}
由 \eqref{eq:1.79} 显然可见 $\mathrm{e}^{\lambda t/(2s)}g(t)$ 单调非降, 于是
\eqref{eq:1.76} 立即成立.
\end{proof}

若 $\Delta_{\Sigma}f\geq0$, 则称 $f$ 在 $\Sigma$ 上是
\newterm{subharmonic}{次调和}{subharmonic}的. 对非负次调和函数这一特殊
情形, 我们立即得到下述平均值不等式:

\begin{corollary}
\label{cor:1.17}
设 $\Sigma^k\subset\R^n$ 是极小子流形, $x_0\in\R^n$, $f$ 是 $\Sigma$ 上的
非负次调和函数; 则
\begin{equation}
s^{-k}\int_{B_s(x_0)\cap\Sigma}f
\label{eq:1.80}
\end{equation}
是 $s$ 的非降函数. 特别地, 若 $x_0\in\Sigma$, 则对一切 $s>0$,
\begin{equation}
f(x_0)\leq\frac{\int_{B_s(x_0)\cap\Sigma}f}{\Vol(B_s\subset\R^k)}.
\label{eq:1.81}
\end{equation}
\end{corollary}

在第 3 章中, 当 $\Sigma$ 不是极小的时候, 我们需要一个一般的平均值不等式.
我们把它记在下面:

\begin{lemma}
\label{lem:1.18}
若 $\Sigma^n\subset\R^{n+1}$ 是平均曲率为 $H$ 的超曲面, $f$ 是 $\Sigma$ 上
的非负函数, 则对 $s<t$,
\begin{equation}
\begin{aligned}
&t^{-k}\int_{B_t\cap\Sigma}f-s^{-k}\int_{B_s\cap\Sigma}f\\
&\quad\geq\int_s^t\tau^{-n-1}\int_{B_\tau\cap\Sigma}
 \langle x,\nabla f+fHN\rangle\,d\tau.
\end{aligned}
\label{eq:1.82}
\end{equation}
\end{lemma}

\begin{proof}
证明与命题~\ref{prop:1.15} 相同, 只是要保留含 $H$ 的项, 并舍去含 $x^{N}$
的非负(因而只是有利的)项. 细节留作练习.
\end{proof}

% ================================================================
\section{高斯映射}
\label{sec:1.4}

设 $\Sigma^2\subset\R^3$ 是曲面. 高斯映射是单位法向量的一个连续选取
\begin{equation}
N:\Sigma\to\bS^2\subset\R^3.
\label{eq:1.83}
\end{equation}
注意这样的映射有两种选择, 即 $N$ 与 $-N$, 它们对应于 $\Sigma$ 的定向的选取.
此外, 映射 $N$ 的微分可以与前面定义的 Weingarten 映射等同起来. 为看出
这一点, 设 $E_1,E_2$ 是 $\Sigma$ 上的正交标架. 由于 $\bS^2$ 在 $N(x)$ 处的
单位法向量恰是 $N(x)$ 本身, $E_1,E_2$ 也就给出像集上的正交标架. 利用这一
等同, 微分 $dN$ 由
\begin{equation}
\langle dN(E_i),E_j\rangle
 =\langle\nabla_{E_i}N,E_j\rangle
 =-\langle N,\nabla_{E_i}E_j\rangle=-A_{ij}
\label{eq:1.84}
\end{equation}
给出. 在最后一个等式中, 我们把法向量 $A_{ij}$ 与它和 $N$ 的内积等同了
起来(因为 $\Sigma$ 是超曲面).

若 $\Sigma$ 极小, 则由于 Weingarten 映射的特征值为 $\kappa_1$ 与
$\kappa_2=-\kappa_1$, 高斯映射是(反)共形的. 此外, 对极小曲面有
\begin{equation}
|dN|^2=|A|^2=\kappa_1^2+\kappa_2^2
 =-2\kappa_1\kappa_2=-2K=-2\det(dN),
\label{eq:1.85}
\end{equation}
而高斯映射的面积是总曲率的倍数. $\R^3$ 中极小曲面的高斯映射的这种共形性,
即 \eqref{eq:1.85}, 可用来证明下一节要讲的经典 Bernstein 定理.

\mtpara{图上的局部坐标}
我们以计算图的度量、曲率与第二基本形式来结束本节. 若 $\Sigma\subset\R^3$
是函数 $u=u(x,y)$ 的图, 则如我们已经看到的, 可以取
\begin{equation}
N=\frac{(-u_x,-u_y,1)}{\sqrt{1+|\nabla u|^2}}.
\label{eq:1.86}
\end{equation}
以 $(x,y)$ 作为图上的坐标, 诱导度量 $g$ 可以表示为
\begin{equation}
g_{xx}=\bigl(1+u_x^2\bigr),\quad
g_{xy}=g_{yx}=u_xu_y,\quad
g_{yy}=\bigl(1+u_y^2\bigr).
\label{eq:1.87}
\end{equation}
直接计算可知, 矩阵 $g$ 的特征值为 $1$ 与 $(1+|\nabla u|^2)$. 这一点在几何上
也很容易看出. 类似地, 逆矩阵为
\begin{equation}
g^{xx}=\frac{1+u_y^2}{1+|\nabla u|^2},\quad
g^{xy}=g^{yx}=\frac{-u_xu_y}{1+|\nabla u|^2},\quad
g^{yy}=\frac{1+u_x^2}{1+|\nabla u|^2}.
\label{eq:1.88}
\end{equation}
由 Gauss 方程, $u$ 的图的 Gauss 曲率为
\begin{equation}
\begin{aligned}
K&=\kappa_1\kappa_2
 =\frac{\langle N_x,(1,0,u_x)\rangle\langle N_y,(0,1,u_y)\rangle
  -\langle N_x,(0,1,u_y)\rangle^{2}}
 {\bigl|(1,0,u_x)\times(0,1,u_y)\bigr|^{2}}\\
&=\frac{\langle(-u_{xx},-u_{yx},0),(1,0,u_x)\rangle
  \langle(-u_{xy},-u_{yy},0),(0,1,u_y)\rangle}
 {(1+|\nabla u|^2)^{2}}\\
&\quad-\frac{\langle(-u_{xx},-u_{xy},0),(0,1,u_y)\rangle^{2}}
 {(1+|\nabla u|^2)^{2}}
 =\frac{u_{xx}u_{yy}-u_{xy}^2}{(1+|\nabla u|^2)^{2}}.
\end{aligned}
\label{eq:1.89}
\end{equation}
因此,
\begin{equation}
K\,d\Area
 =\frac{u_{xx}u_{yy}-u_{xy}^2}{(1+|\nabla u|^2)^{3/2}}\,dx\wedge dy.
\label{eq:1.90}
\end{equation}
类似地, 我们可以在坐标 $(x,y)$ 下把第二基本形式 $A$ 表示为
\begin{equation}
A_{xx}=\bigl\langle\nabla_{(1,0,u_x)}(1,0,u_x),N\bigr\rangle
 =\bigl\langle(0,0,u_{xx}),N\bigr\rangle
 =\frac{u_{xx}}{(1+|\nabla u|^2)^{1/2}}.
\label{eq:1.91}
\end{equation}
于是, 这一计算(以及类似的计算)给出
\begin{equation}
\begin{aligned}
A_{xx}&=\frac{u_{xx}}{(1+|\nabla u|^2)^{1/2}},\\
A_{xy}&=A_{yx}=\frac{u_{xy}}{(1+|\nabla u|^2)^{1/2}},\\
A_{yy}&=\frac{u_{yy}}{(1+|\nabla u|^2)^{1/2}}.
\end{aligned}
\label{eq:1.92}
\end{equation}

回忆(由 \eqref{eq:1.27}), $A$ 的范数平方是 $A$ 在正交标架下各分量的平方和; 在一般
标架下, 它由
\begin{equation}
|A|^2=A_{ij}A_{kl}g^{ik}g^{jl}
\label{eq:1.93}
\end{equation}
给出. 第二基本形式的表达式与 $g$ 的特征值估计(见 \eqref{eq:1.87})合起来蕴含
\begin{equation}
\frac{\bigl|\operatorname{Hess}_u\bigr|^{2}}{(1+|\nabla u|^2)^{3}}
 \leq|A|^{2}\leq
 2\frac{\bigl|\operatorname{Hess}_u\bigr|^{2}}{1+|\nabla u|^2}.
\label{eq:1.94}
\end{equation}

% ================================================================
\section{Bernstein 定理}
\label{sec:1.5}

在证明著名的 Bernstein 定理之前, 我们先给出极小图总曲率的一个界. 后面
将会看到, 再做一些工作, 这个界就可以用来给出极小图的局部曲率估计. 这样的
局部曲率估计最初由 Heinz \cite{He} 用复分析证明, 它给出了 Bernstein 定理
的一个推广.

\begin{lemma}
\label{lem:1.19}
若 $u:\Omega\subset\R^2\to\R$ 是极小曲面方程的解, 则对一切支撑含于
$\Omega\times\R$ 中的非负 Lipschitz 函数 $\eta$,
\begin{equation}
\int_{\Graph_u}|A|^{2}\eta^{2}
 \leq C\int_{\Graph_u}\bigl|\nabla_{\Graph_u}\eta\bigr|^{2}.
\label{eq:1.95}
\end{equation}
\end{lemma}

\begin{proof}
设 $\omega$ 表示单位球面 $\bS^2$ 上的面积 $2$-形式. 由于上半球是可缩的,
可以找到上半球上的一个(光滑)$1$-形式 $\alpha$, 使得
\begin{equation}
d\alpha=\omega.
\label{eq:1.96}
\end{equation}
如前, 用 $N$ 表示高斯映射. 由于 $\Sigma$ 极小, 且微分 $d$ 与拉回可交换,
我们看到
\begin{equation}
|A|^{2}\,d\Area=-2K\,d\Area=2N^{*}\omega=2\,dN^{*}\alpha.
\label{eq:1.97}
\end{equation}
此外, 由于 $\alpha$ 是 $1$-形式, 存在常数 $C_{\alpha}$ 使得
\begin{equation}
|N^{*}\alpha|\leq C_{\alpha}|dN|=C_{\alpha}|A|.
\label{eq:1.98}
\end{equation}
令 $\Sigma=\Graph_u$. 由 \eqref{eq:1.97}、Stokes 定理与 \eqref{eq:1.98}, 得到
\begin{equation}
\begin{aligned}
\int_{\Sigma}\eta^{2}|A|^{2}\,d\Area
&=2\int_{\Sigma}\eta^{2}\,dN^{*}\alpha
 =-4\int_{\Sigma}\eta\,d\eta\wedge N^{*}\alpha\\
&\leq4C_{\alpha}\int_{\Sigma}\eta\,\bigl|\nabla_{\Sigma}\eta\bigr|
 |A|\,d\Area\\
&\leq4C_{\alpha}
 \Bigl(\int_{\Sigma}\eta^{2}|A|^{2}\,d\Area\Bigr)^{\frac12}
 \Bigl(\int_{\Sigma}\bigl|\nabla_{\Sigma}\eta\bigr|^{2}\,d\Area\Bigr)^{\frac12},
\end{aligned}
\label{eq:1.99}
\end{equation}
其中最后一个不等式用了 Cauchy--Schwarz 不等式. 最后, 把它平方并两边除以
$\int_{\Sigma}\eta^{2}|A|^{2}$, 得
\begin{equation}
\int_{\Sigma}\eta^{2}|A|^{2}
 \leq16C_{\alpha}^{2}\int_{\Sigma}
 \bigl|\nabla_{\Sigma}\eta\bigr|^{2}.
\label{eq:1.100}
\end{equation}
这就证明了引理.
\end{proof}

\mtpara{对数截断技巧}
上一个引理表明, 可以用截断函数 $\eta$ 的能量来控制总曲率. 当图 $\Sigma$ 是
整的(即定义在整个 $\R^2$ 上)时, 我们将看到可以选取一列 $\eta$ 收敛到 $1$,
而其能量趋于零. 当 $\Sigma$ 具有二次面积增长时, 这是可以做到的, 这一技巧
称为``对数截断技巧''.

为看出这一点, 固定一个大整数 $N$, 在 $B_{\mathrm{e}^{2N}}\subset\R^2$ 上
定义 $\eta$ 为
\begin{equation}
\eta=
\begin{cases}
1, & \text{若 } r\leq\mathrm{e}^{N},\\
2-\log(r)/N, & \text{若 } \mathrm{e}^{N}<r\leq\mathrm{e}^{2N},\\
0, & \text{若 } \mathrm{e}^{2N}<r,
\end{cases}
\label{eq:1.101}
\end{equation}
其中 $r=|x|$. 由于 $|\nabla\eta|=\frac{1}{Nr}$ 是径向的, $\eta$ 的能量为
\begin{equation}
\begin{aligned}
\int_{\R^2}|\nabla\eta|^{2}
&=2\pi\int_{\mathrm{e}^{N}}^{\mathrm{e}^{2N}}\frac{1}{N^{2}r^{2}}\,r\,dr
 =\frac{2\pi}{N^{2}}\int_{\mathrm{e}^{N}}^{\mathrm{e}^{2N}}\frac{dr}{r}\\
&=\frac{2\pi}{N^{2}}\bigl[\log(\mathrm{e}^{2N})-\log(\mathrm{e}^{N})\bigr]
 =\frac{2\pi}{N}.
\end{aligned}
\label{eq:1.102}
\end{equation}
特别地, 取一列趋于无穷的 $N$, 便得到一列收敛到 $1$ 且能量趋于 $0$ 的截断
函数. 显然, 同样的论证适用于把 $r$ 取为流形 $\Sigma$ 上到某个固定点的距离,
只要对某个固定常数 $C$ 有
\[
\Vol(\partial B_r)\leq Cr.
\]
不过, 当 $\Sigma$ 具有二次面积增长
\[
\Vol(B_r)\leq Cr^{2}
\]
时, 只要把积分拆成内圆与外圆半径之比有界的圆环上的积分, 论证仍可走通. 即:
利用
\[
\sup_{B_{\mathrm{e}^{\ell}}\setminus B_{\mathrm{e}^{\ell-1}}}
|\nabla\eta|^{2}\leq N^{-2}\mathrm{e}^{2-2\ell},
\]
得到
\begin{equation}
\begin{aligned}
\int|\nabla\eta|^{2}
&\leq\sum_{\ell=N+1}^{2N}
 \int_{B_{\mathrm{e}^{\ell}}\setminus B_{\mathrm{e}^{\ell-1}}}
 \Bigl[N^{-2}\mathrm{e}^{2-2\ell}\Vol\bigl(B_{\mathrm{e}^{\ell}}\bigr)\Bigr]\\
&\leq C\mathrm{e}^{2}N^{-2}\sum_{\ell=N+1}^{2N}
 \int_{B_{\mathrm{e}^{\ell}}\setminus B_{\mathrm{e}^{\ell-1}}}
 =\frac{C\mathrm{e}^{2}}{N}.
\end{aligned}
\label{eq:1.103}
\end{equation}

接下来, 我们把上一个引理与对数截断技巧结合起来, 得到极小图总曲率的一个界.

\begin{corollary}
\label{cor:1.20}
若 $u:\Omega\subset\R^2\to\R$ 是极小曲面方程的解, $\kappa>1$, 且 $\Omega$
包含以原点为心、半径为 $\kappa R$ 的球, 则
\begin{equation}
\int_{B_{\sqrt{\kappa}R}\cap\Graph_u}|A|^{2}
 \leq\frac{C}{\log\kappa}.
\label{eq:1.104}
\end{equation}
\end{corollary}

\begin{proof}
令 $\Sigma=\Graph_u$. 在整个 $\R^3$ 上定义截断函数 $\eta$, 然后把它限制到
$u$ 的图上, 做法如下: 用 $r$ 表示 $\R^3$ 中到原点的距离, 定义 $\eta$ 为
\begin{equation}
\eta=
\begin{cases}
1, & \text{若 } r^{2}\leq\kappa R^{2},\\
2-2\log(rR^{-1})/\log\kappa,
  & \text{若 } \kappa R^{2}<r^{2}\leq\kappa^{2}R^{2},\\
0, & \text{若 } r^{2}>\kappa^{2}R^{2}.
\end{cases}
\label{eq:1.105}
\end{equation}
由于 $|\nabla_{\Sigma}r|\leq|\nabla r|=1$, 我们有
\begin{equation}
\bigl|\nabla_{\Sigma}\eta\bigr|
 \leq\frac{2}{r\log\kappa}.
\label{eq:1.106}
\end{equation}
用这个截断函数 $\eta$ 应用引理~\ref{lem:1.19}, 并利用面积界 \eqref{eq:1.20}, 得
\begin{equation}
\begin{aligned}
\int_{B_{\sqrt{\kappa}R}\cap\Sigma}|A|^{2}
&\leq\int_{\Sigma}\eta^{2}|A|^{2}
 \leq C\int_{\Sigma}\bigl|\nabla_{\Sigma}\eta\bigr|^{2}\\
&\leq\frac{4C}{(\log\kappa)^{2}}
 \int_{B_{\kappa R}\cap\Sigma}r^{-2}\\
&\leq\frac{4C}{(\log\kappa)^{2}}
 \sum_{\ell=(\log\kappa)/2}^{\log\kappa}
 \int_{(B_{\mathrm{e}^{\ell}R}\setminus B_{\mathrm{e}^{\ell-1}R})\cap\Sigma}
 r^{-2}\\
&\leq\frac{4C}{(\log\kappa)^{2}}
 \sum_{\ell=(\log\kappa)/2}^{\log\kappa}2\pi\mathrm{e}^{2}
 \leq\frac{4\pi C\mathrm{e}^{2}}{\log\kappa}.
\end{aligned}
\label{eq:1.107}
\end{equation}
(这里为简单起见, 假定了 $(\log\kappa)/2$ 是整数.)
\end{proof}

如上所述, 这一论证, 即取上述特殊的 $\eta$ 作分部积分, 通常称为``对数截断
论证''. 当曲面至多具有二次面积增长时, 它相当有用.

作为这个推论的结论, 我们得到 S. Bernstein \cite{Be} 1916 年的下述定理:

\begin{theorem}[Bernstein, \cite{Be}]
\label{thm:1.21}
若 $u:\R^2\to\R$ 是极小曲面方程的整解, 则 $u(x,y)=ax+by+c$, 其中
$a,b,c\in\R$ 为常数.
\end{theorem}

\begin{proof}
由上一个推论, 对一切 $R>1$ 有
\begin{equation}
\int_{B_{\sqrt{R}}\cap\Graph_u}|A|^{2}\leq\frac{C}{\log R}.
\label{eq:1.108}
\end{equation}
令 $R\to\infty$, 我们断定 $|A|^{2}\equiv0$; 于是
$0=u_{xx}=u_{xy}=u_{yy}$, 从而 $u=ax+by+c$, 其中 $a,b,c\in\R$ 为常数.
\end{proof}

上述证明(来自 L. Simon \cite{Si5})在两个方面用到了极小性:

\begin{itemize}
  \item 极小图的面积界, 即 \eqref{eq:1.20}.
  \item 高斯映射的共形性, 即 \eqref{eq:1.85}.
\end{itemize}

这个证明其实可以用于更广的一类微分方程, 只要把共形性换成拟共形性. 我们
以后会简要回到这一点(在 \eqref{eq:7.7} 处, 那里也会定义拟共形性), 但在本讲义中不
讨论拟共形映射的估计. 详细的讨论可见 \cite{GiTr} 的第 16 章.

此外, 只要高斯映射漏掉 $\bS^2$ 中的一个开集, 引理~\ref{lem:1.19} 的论证
就适用. 利用这一点, 便不难把 Bernstein 定理推广到高斯映射漏掉一个开集的
完全极小曲面; 这是 L. Nirenberg 的一个猜想, 由 R. Osserman 用另一种方法
证明, 见 \cite{Os4}.

% ================================================================
\section{Weierstrass 表示}
\label{sec:1.6}

经典的 Weierstrass 表示取全纯数据(一个 Riemann 曲面、一个亚纯函数与一个
全纯 $1$-形式), 并由此得到 $\R^3$ 中的一张极小曲面. 确切地说, 给定
\begin{itemize}
  \item 一个 Riemann 曲面 $\Omega$,
  \item $\Omega$ 上的亚纯函数 $g$,
  \item $\Omega$ 上的全纯 $1$-形式 $\phi$,
\end{itemize}
则由
\begin{equation}
F(z)=\operatorname{Re}\int_{\zeta\in\gamma_{z_0,z}}
\Bigl(\frac12\bigl(g^{-1}(\zeta)-g(\zeta)\bigr),\,
      \frac{\ii}{2}\bigl(g^{-1}(\zeta)+g(\zeta)\bigr),\,1\Bigr)\phi(\zeta)
\label{eq:1.109}
\end{equation}
得到(分支)共形极小浸入 $F:\Omega\to\R^3$. 这里 $z_0\in\Omega$ 是固定的基点,
积分沿从 $z_0$ 到 $z$ 的路径 $\gamma_{z_0,z}$ 进行. $z_0$ 的选取只使 $F$
相差一个常数. 一般说来, 映射 $F$ 可能依赖于路径的选取(因而未必定义良好);
这称为``周期问题''(最新的进展见 M. Weber 与 M. Wolf, \cite{WeWo2}).

\begin{lemma}
\label{lem:1.22}
若 $f^1,f^2,f^3$ 是 $\Omega\subset\mathbf{C}$ 上的全纯函数, 且
$F=(F^1,F^2,F^3):\Omega\to\R^3$ 由
\begin{equation}
F(z)=\operatorname{Re}\int_{\zeta\in\gamma_{z_0,z}}
\bigl(f^1,f^2,f^3\bigr)\,d\zeta
\label{eq:1.110}
\end{equation}
给出, 则对每个 $i=1,2,3$ 都有
\begin{equation}
\frac{\partial F^i}{\partial x}
 -\ii\frac{\partial F^i}{\partial y}=f^i.
\label{eq:1.111}
\end{equation}
\end{lemma}

\begin{proof}
若 $g=u+\ii v$ 全纯, 则 Cauchy--Riemann 方程为
\begin{equation}
\frac{\partial u}{\partial x}=\frac{\partial v}{\partial y},
\label{eq:1.112}
\end{equation}
\begin{equation}
\frac{\partial u}{\partial y}=-\frac{\partial v}{\partial x}.
\label{eq:1.113}
\end{equation}
特别地, 我们得到
\begin{equation}
\begin{aligned}
2g'&=\frac{\partial g}{\partial x}-\ii\frac{\partial g}{\partial y}
 =\Bigl(\frac{\partial u}{\partial x}+\frac{\partial v}{\partial y}\Bigr)
 -\ii\Bigl(\frac{\partial u}{\partial y}-\frac{\partial v}{\partial x}\Bigr)\\
&=2\Bigl(\frac{\partial u}{\partial x}
 -\ii\frac{\partial u}{\partial y}\Bigr).
\end{aligned}
\label{eq:1.114}
\end{equation}
\end{proof}

{\setlength{\emergencystretch}{2em}%
作为引理~\ref{lem:1.22} 的直接推论, 我们看到 Weierstrass 表示给出的是共形
映射(见下面的 \eqref{eq:1.116}). 由于在二维情形调和函数是共形不变的, 坐标函数限制在
Weierstrass 表示的像上便是调和函数. 于是我们看到, 像事实上是极小曲面.\par}

由此得到的曲面的单位法向量 $N$、度量 $ds^2$ 与 Gauss 曲率 $K$ 分别为
\begin{equation}
N=\frac{(2\operatorname{Re}g,\,2\operatorname{Im}g,\,|g|^2-1)}{|g|^2+1},
\label{eq:1.115}
\end{equation}
\begin{equation}
ds^2=\frac{|\phi|^2}{4}\bigl(|g|^{-1}+|g|\bigr)^2,
\label{eq:1.116}
\end{equation}
\begin{equation}
K=-\Bigl[\frac{4|\partial_z g|\,|g|}{|\phi|\,(1+|g|^2)^2}\Bigr]^2.
\label{eq:1.117}
\end{equation}

由于由 \eqref{eq:1.109}, 拉回 $F^{*}(dx_3)$ 是 $\operatorname{Re}\phi$, 故 $\phi$
通常称为\emph{高度微分}. 由 \eqref{eq:1.115}, $g$ 是高斯映射再作球极投影所得的复合.

由 Weierstrass 数据构造极小曲面的标准例子是
\begin{equation}
g(z)=z,\quad\phi(z)=dz/z,\quad
\Omega=\mathbf{C}\setminus\{0\}\ \text{给出悬链面},
\label{eq:1.118}
\end{equation}
\begin{equation}
g(z)=\mathrm{e}^{\ii z},\quad\phi(z)=dz,\quad
\Omega=\mathbf{C}\ \text{给出螺旋面},
\label{eq:1.119}
\end{equation}
\begin{equation}
g(z)=1/z,\quad\phi(z)=2z\,dz,\quad
\Omega=\mathbf{C}\ \text{给出 Enneper 曲面}.
\label{eq:1.120}
\end{equation}

表示 \eqref{eq:1.109} 给出了形变极小浸入的一个显然途径. 即: 把 $1$-形式 $\phi$ 乘上
一个单位复数 $\mathrm{e}^{\ii\theta}$, 便得到另一个极小浸入. 此外, 由于度量
$ds^2$ 只依赖于 $|\phi|$, 这张新曲面与原曲面等距. 当 $\theta=\pi/2$ 时, 新
曲面称为\emph{共轭极小曲面}; 对一般的 $\theta$, 它们称为\emph{伴随极小曲面}.
通过变量替换不难看出, 螺旋面与(万有覆盖的)悬链面是共轭极小曲面.

Weierstrass 表示在构造浸入极小曲面时特别有用. 例如, 在 \cite{Na1} 中,
Nadirashvili 用它构造了 $\R^3$ 单位球中的一个完全浸入极小曲面. 特别地,
Nadirashvili 的曲面不是固有的, 即它与紧集的交未必紧.

通常, 证明所得的浸入是嵌入(即一一的)相当困难, 尽管有一些有趣的情形可以做
到. 第一个现代例子是 \cite{HoMe1}, 其中 D. Hoffman 与 W. Meeks 证明了 Costa
构造的曲面是嵌入的; 这是自经典的悬链面、螺旋面与平面以来所发现的第一个新的
完全、有限拓扑、真嵌入极小曲面. 这引出了许多这类曲面的发现(进一步的讨论见
\cite{HoK}, \cite{Os2} 与 \cite{Ro}).

\mtpara{有限总曲率}
若
\[
\int_{\Sigma}|A|^{2}<\infty,
\]
则称曲面 $\Sigma^2\subset\R^n$ 具有\emph{有限总曲率}. 当 $\Sigma\subset\R^3$
极小的时候, Gauss 方程蕴含这等价于
\[
\int_{\Sigma}|K_{\Sigma}|=-\int_{\Sigma}K_{\Sigma}<\infty.
\]
1964 年, R. Osserman \cite{Os3} 证明了: $\R^3$ 中的完全极小曲面共形等价于
一个闭 Riemann 曲面去掉有限多个点, 且高斯映射可以亚纯地延拓过这些穿孔点.

\mtpara{Schwarz 反射原理}
Schwarz 反射原理给出一种把微分方程的解延拓过边界的办法, 前提是解沿边界具有
较好的形状. 我们将看到它对调和函数、全纯函数与极小曲面的几种版本.

我们从调和函数的反射原理开始. 第一个观察是: 调和函数继承其边界数据的
对称性:

\begin{lemma}
\label{lem:1.23}
设 $v:D\to\R$ 在 $D$ 上调和、在 $\overline{D}$ 上连续.
\begin{enumerate}
  \item 若在 $\partial D$ 上 $v(x,y)=v(x,-y)$, 则在 $\overline{D}$ 上
    $v(x,y)=v(x,-y)$.
  \item 若在 $\partial D$ 上 $v(x,y)=-v(x,-y)$, 则在 $\overline{D}$ 上
    $v(x,y)=-v(x,-y)$.
\end{enumerate}
\end{lemma}

\begin{proof}
定义反射函数 $w:D\to\R$ 为 $w(x,y)=v(x,-y)$. 由链式法则, $w$ 也是调和
函数. 在第一种情形, $w$ 与 $v$ 在 $\partial D$ 上相等, 故由极大值原理,
$w=v$ 于 $\overline{D}$. 第二种情形中同样的论证给出 $w=-v$ 于
$\overline{D}$.
\end{proof}

设 $D^+=D_1\cap\{y>0\}$ 表示上半单位圆盘. 现在我们可以证明调和函数的反射
原理.

\begin{lemma}
\label{lem:1.24}
若 $u:D^+\to\R$ 调和且 $u(x,0)=0$, 则按
\begin{equation}
u(x,y)=-u(x,-y)
\label{eq:1.121}
\end{equation}
把 $u$ 延拓到整个 $D$, 得到 $D$ 上的调和函数.
\end{lemma}

\begin{proof}
用 \eqref{eq:1.121} 在 $D$ 上定义 $u$, 并设 $v:D\to\R$ 是在边界上与 $u$ 相等的调和
函数\footnote{这个函数由解 Dirichlet 问题给出; 见 \cite{GiTr}、\cite{A} 或
第 4 章.}. 我们将证明 $v=u$ 于 $D$.

注意 $u$ 的构造保证在 $\partial D$ 上 $w(x,y)=-w(x,-y)$, 从而由
引理~\ref{lem:1.23} 的 (2), 在整个 $D$ 上 $w(x,y)=-w(x,-y)$. 特别地, 这
给出
\[
w(x,0)=0.
\]
因此 $w$ 与 $u$ 是 $D^+$ 上的调和函数, 且在 $\partial D^+$ 上相等(它们在
$\partial D^+\cap\partial D$ 上显然相等, 而在 $\partial D^+$ 的其余部分上
都为零). 由极大值原理, $w=u$ 于 $\overline{D^+}$. 同样的论证适用于
$D\setminus D^+$, 引理得证.
\end{proof}

同样的论证(见 \cite{A} 第 172 页)表明: 若 $f$ 在 $D^+$ 上全纯, 且在
$\{y=0\}$ 上只取实值, 则可用
\[
f(\bar z)=\overline{f(z)}
\]
把它延拓为整个 $D$ 上的全纯函数.

Schwarz 为极小曲面发展了反射原理的一个版本, 并用它构造三周期极小曲面.
我们陈述 Lawson \cite{La2} 的一个版本, 证明可参看该处:

\begin{proposition}[Schwarz 反射]
\label{prop:1.25}
设 $u:D^+\to\R^N$ 是极小曲面, 且 $u$ 把 $\{y=0\}$ 映到一条直线 $\gamma$
上. 则关于 $\gamma$ 的反射给出极小曲面 $u:D\to\R^N$.
\end{proposition}

% ================================================================
\section{强极值原理}
\label{sec:1.7}

首先注意, 极小曲面方程两个解之差满足一个一致椭圆的散度型方程(其中椭圆性的
界依赖于极小图的梯度的界):

\begin{lemma}
\label{lem:1.26}
若 $u_1$ 与 $u_2$ 是区域 $\Omega\subset\R^n$ 上极小曲面方程的解, 则
$v=u_2-u_1$ 满足形如
\begin{equation}
\dive\bigl(a_{i,j}\nabla v\bigr)=0
\label{eq:1.122}
\end{equation}
的方程, 其中矩阵 $a_{i,j}=a_{i,j}(x)$ 的特征值满足
\begin{equation}
0<\mu\leq\lambda_1\leq\dots\leq\lambda_n\leq1/\mu.
\label{eq:1.123}
\end{equation}
常数 $\mu$ 只依赖于 $|\nabla u_i|$ 的上界.
\end{lemma}

\begin{proof}
定义映射 $F:\R^n\to\R^n$ 为
\begin{equation}
F(X)=\frac{X}{(1+|X|^2)^{1/2}}.
\label{eq:1.124}
\end{equation}
由微积分基本定理与链式法则, 可以写成
\begin{equation}
\begin{aligned}
F(\nabla u_2)-F(\nabla u_1)
&=\int_0^1\frac{d}{dt}
 \Bigl(F\bigl(\nabla u_1+t(\nabla u_2-\nabla u_1)\bigr)\Bigr)dt\\
&=\int_0^1 dF\bigl(\nabla u_1+t(\nabla u_2-\nabla u_1)\bigr)
 \nabla(u_2-u_1)\,dt\\
&=\Bigl(\int_0^1 dF\bigl(\nabla u_1+t(\nabla u_2-\nabla u_1)\bigr)dt
 \Bigr)\nabla(u_2-u_1).
\end{aligned}
\label{eq:1.125}
\end{equation}
由此我们断定 $v=u_2-u_1$ 满足形如
\begin{equation}
\dive\bigl(a_{i,j}\nabla v\bigr)=0
\label{eq:1.126}
\end{equation}
的方程, 其中矩阵 $a_{i,j}$ 由 \eqref{eq:1.125} 给出.

给定单位向量 $V\in\bS^{n-1}$ 与 $X\in\R^n$, 我们看到
\begin{equation}
dF(X)V=\frac{V}{(1+|X|^2)^{1/2}}
 -\frac{\langle X,V\rangle}{(1+|X|^2)^{3/2}}X.
\label{eq:1.127}
\end{equation}
特别地, 与 $V$ 取内积给出
\begin{equation}
\begin{aligned}
(1+|X|^2)^{3/2}\bigl\langle V,dF(X)V\bigr\rangle
&=(1+|X|^2)-\langle X,V\rangle^{2}\\
&\geq(1+|X|^2)-|X|^{2}=1.
\end{aligned}
\label{eq:1.128}
\end{equation}
于是 $a_{i,j}$ 是正定矩阵的加权平均, 因而也正定, 证明完成.
\end{proof}

下述推论是极小超曲面的强极值原理的局部形式:

\begin{corollary}
\label{cor:1.27}
设 $\Omega\subset\R^n$ 是原点的开连通邻域. 若 $u_1,u_2:\Omega\to\R$ 是
极小曲面方程的解, 满足 $u_1\leq u_2$ 且 $u_1(0)=u_2(0)$, 则
$u_1\equiv u_2$.
\end{corollary}

\begin{proof}
由于引理~\ref{lem:1.26} 给出的矩阵 $a_{i,j}$ 正定, 我们可以对 $v=u_2-u_1$
应用线性方程的极大值原理(例如见 \cite{HaLi}, 或 \cite{GiTr} 的定理 3.5).
\end{proof}

把超曲面局部写成函数的图, 便看出推论~\ref{cor:1.27} 有下述直接推论:

\begin{corollary}[强极值原理]
\label{cor:1.28}
若 $\Sigma_1,\Sigma_2\subset\R^n$ 是完全连通的极小超曲面(无边界),
$\Sigma_1\cap\Sigma_2\neq\emptyset$, 且 $\Sigma_2$ 位于 $\Sigma_1$ 的一侧,
则 $\Sigma_1=\Sigma_2$.
\end{corollary}

作为强极值原理的应用, 我们下面证明 Rado 的一个著名定理的 Schoen 推广
(见\thmref{thm:6.14}).

\begin{theorem}[Schoen, \cite{Sc2}]
\label{thm:1.29}
设 $\Omega\subset\R^2$ 严格凸, $\sigma\subset\R^3$ 是一条简单闭曲线, 它为
$\partial\Omega$ 上具有有界斜率的图. 则任何满足 $\partial\Sigma=\sigma$ 的极小曲面
$\Sigma\subset\R^3$ 必定是 $\Omega$ 上的图, 从而唯一.
\end{theorem}

\begin{proof}
我们证明 $\Sigma$ 是图, 把唯一性留给读者. 由极大值原理, $\Sigma$ 的内部
含于柱体 $\Omega\times\R$ 的内部. 给定平面 $\{x_3=t\}$, 我们把 $\Sigma$ 分成
位于该平面上方的部分 $\Sigma_t^+$ 与下方的部分 $\Sigma_t^-$. 把
$\Sigma_t^+$ 关于该平面反射, 得到位于平面下方的新的极小曲面
$\tilde\Sigma_t^+$. 由极大值原理, 不可能存在第一个 $t$, 使 $\Sigma_t^-$ 与
$\tilde\Sigma_t^+$ 有内点接触. 又因 $\partial\Sigma$ 是图, 也不可能有边界
接触点. 于是 $\Sigma$ 到平面 $\{x_3=0\}$ 的投影必是一一的, 正如所愿.
\end{proof}

上面把极大值原理与反射结合起来的做法称为``移动平面法''. 它最初由 Alexandrov
发展, 用来证明 $\R^n$ 中闭的嵌入常平均曲率超曲面必是圆球面.

% ================================================================
\section{第二变分公式、Morse 指标与稳定性}
\label{sec:1.8}

\mtpara{第二变分公式}
现在设 $\Sigma^k\subset M^n$ 是极小子流形; 我们要计算 $\Sigma$ 的一个变分下面积
泛函的二阶导数. 因此, 仍设 $F$ 是 $\Sigma$ 的一个具紧支集的变分. 事实上, 我们
将假定 $F$ 是法向变分, 即在 $\Sigma$ 上有
\begin{equation}
F_t^{T}\equiv0.
\label{eq:1.129}
\end{equation}
如前, 设 $x_i$ 是 $\Sigma$ 上的局部坐标, 并令
\begin{equation}
g_{ij}(t)=g\bigl(F_{x_i},F_{x_j}\bigr),
\label{eq:1.130}
\end{equation}
\begin{equation}
\nu(t)=\sqrt{\det\bigl(g_{ij}(t)\bigr)}\,
 \sqrt{\det\bigl(g^{ij}(0)\bigr)}.
\label{eq:1.131}
\end{equation}
对测度 $\nu(t)$ 求导给出
\begin{equation}
\frac{d^{2}}{dt^{2}}\Big|_{t=0}\Vol\bigl(F(\Sigma,t)\bigr)
 =\int\frac{d^{2}}{dt^{2}}\Big|_{t=0}\nu(t)\,
 \sqrt{\det\bigl(g_{ij}(0)\bigr)}.
\label{eq:1.132}
\end{equation}
回忆测度 $\nu(t)$ 的一阶导数可以写成
\begin{equation}
2\frac{d}{dt}\nu(t)
 =\operatorname{Tr}\bigl(g'_{ij}(t)\,g^{\ell m}(t)\bigr)\nu(t),
\label{eq:1.133}
\end{equation}
这里的迹指 $\sum_{i,j}g'_{ij}(t)g^{ij}(t)$. 为看出这一点, 回忆
\begin{equation}
\frac{d}{dt}\Big|_{t=0}\det\bigl(\delta_{ij}+ta_{ij}\bigr)
 =\operatorname{Tr}\bigl(a_{ij}\bigr).
\label{eq:1.134}
\end{equation}
为在点 $x\in\Sigma$ 处计算 $d^{2}/dt^{2}|_{t=0}\nu(t)$, 我们可以取坐标系
$x_i$ 使其在 $x$ 处正交. 此时对 \eqref{eq:1.133} 求导, 便在 $x$ 处得到
\begin{equation}
2\frac{d^{2}}{dt^{2}}\Big|_{t=0}\nu(t)
 =\operatorname{Tr}\bigl(g''_{ij}(0)\bigr)
 -\operatorname{Tr}\bigl(g'_{ij}(0)g'_{\ell m}(0)\bigr)
 +\frac12\bigl[\operatorname{Tr}\bigl(g'_{ij}(0)\bigr)\bigr]^{2}.
\label{eq:1.135}
\end{equation}
由于 $\Sigma$ 极小, 我们有 $\operatorname{Tr}(g'(0))=0$, 从而得到
\begin{equation}
2\frac{d^{2}}{dt^{2}}\Big|_{t=0}\nu(t)
 =\operatorname{Tr}\bigl(g''(0)\bigr)
 -\operatorname{Tr}\bigl(g'(0)g'(0)\bigr)
 =\operatorname{Tr}\bigl(g''(0)\bigr)-|g'(0)|^{2},
\label{eq:1.136}
\end{equation}
其中最后一个等式用了对称矩阵 $M$ 的 $\operatorname{Tr}(M^{2})=|M|^{2}$.

\begin{lemma}
\label{lem:1.30}
在点 $x$ 处, 有
\begin{equation}
\bigl|g'(0)\bigr|^{2}
 =4\bigl|\bigl\langle A(\cdot,\cdot),F_t\bigr\rangle\bigr|^{2},
\label{eq:1.137}
\end{equation}
\begin{equation}
\begin{aligned}
\operatorname{Tr}\bigl(g''(0)\bigr)
&=2\bigl|\bigl\langle A(\cdot,\cdot),F_t\bigr\rangle\bigr|^{2}
 +2\bigl|\nabla_{\Sigma}^{N}F_t\bigr|^{2}\\
&\quad+2\operatorname{Tr}
 \bigl\langle \mathrm{R}_M(\cdot,F_t)F_t,\cdot\bigr\rangle
 +2\dive_{\Sigma}(F_{tt}).
\end{aligned}
\label{eq:1.138}
\end{equation}
\end{lemma}

\begin{proof}
一个简单的计算给出
\begin{equation}
g'_{ij}(0)
 =\bigl\langle F_{x_it},F_{x_j}\bigr\rangle
 +\bigl\langle F_{x_i},F_{x_jt}\bigr\rangle
 =-2\bigl\langle A\bigl(F_{x_i},F_{x_j}\bigr),F_t\bigr\rangle.
\label{eq:1.139}
\end{equation}
由于 $g_{ij}$ 在 $x$ 处是单位矩阵, 向量 $F_{x_i}$ 给出 $x$ 处 $T\Sigma$ 的一组
正交基, 于是得到 \eqref{eq:1.137}.

为得到 \eqref{eq:1.138}, 我们计算
\begin{equation}
\operatorname{Tr}g''(0)
 =2\sum_{i=1}^{k}\bigl\langle F_{x_itt},F_{x_i}\bigr\rangle
 +2\sum_{i=1}^{k}\bigl\langle F_{x_it},F_{x_it}\bigr\rangle.
\label{eq:1.140}
\end{equation}
其次, 利用 $M$ 的 Riemann 曲率张量 $R_M$ 的定义\footnote{如 \cite{dC2} 第
89 页那样, 令
$R_M(U,V)W=\nabla_V\nabla_UW-\nabla_U\nabla_VW+\nabla_{[U,V]}W$.}, 得到
\begin{equation}
\begin{aligned}
\sum_{i=1}^{k}\bigl\langle F_{x_itt},F_{x_i}\bigr\rangle
&=\sum_{i=1}^{k}\bigl\langle\nabla_{F_t}\nabla_{F_t}F_{x_i},F_{x_i}\bigr\rangle
 =\sum_{i=1}^{k}\bigl\langle\nabla_{F_t}\nabla_{F_{x_i}}F_t,F_{x_i}\bigr\rangle\\
&=\sum_{i=1}^{k}\bigl\langle R_M(F_{x_i},F_t)F_t,F_{x_i}\bigr\rangle
 +\sum_{i=1}^{k}\bigl\langle\nabla_{F_{x_i}}\nabla_{F_t}F_t,F_{x_i}\bigr\rangle\\
&=\sum_{i=1}^{k}\bigl\langle R_M(F_{x_i},F_t)F_t,F_{x_i}\bigr\rangle
 +\dive_{\Sigma}(F_{tt}).
\end{aligned}
\label{eq:1.141}
\end{equation}
把这些结合起来, 并再次利用 $F_{x_i}$ 是正交标架, 便得到
\begin{equation*}
\begin{aligned}
\operatorname{Tr}g''(0)
&=2\sum_{i=1}^{k}\bigl\langle F_{x_it}^{T},F_{x_it}^{T}\bigr\rangle
 +2\sum_{i=1}^{k}\bigl\langle F_{x_it}^{N},F_{x_it}^{N}\bigr\rangle\\
&\quad+2\sum_{i=1}^{k}\bigl\langle \mathrm{R}_M(F_{x_i},F_t)F_t,F_{x_i}\bigr\rangle
 +2\dive_{\Sigma}(F_{tt})\\
&=2\bigl|\bigl\langle A(\cdot,\cdot),F_t\bigr\rangle\bigr|^{2}
 +2\bigl|\nabla_{\Sigma}^{N}F_t\bigr|^{2}
 +2\operatorname{Tr}_{\Sigma}
 \bigl\langle \mathrm{R}_M(\cdot,F_t)F_t,\cdot\bigr\rangle
 +2\dive_{\Sigma}(F_{tt}).
\end{aligned}
\end{equation*}
因此, 在 $x$ 处得到
\begin{equation}
\begin{aligned}
\frac{d^{2}}{dt^{2}}\Big|_{t=0}\nu(t)
&=-\bigl|\bigl\langle A(\cdot,\cdot),F_t\bigr\rangle\bigr|^{2}\\
&\quad+\bigl|\nabla_{\Sigma}^{N}F_t\bigr|^{2}
 -\operatorname{Tr}_{\Sigma}
 \bigl\langle \mathrm{R}_M(E_i,F_t)E_i,F_t\bigr\rangle
 +\dive_{\Sigma}(F_{tt}).
\end{aligned}
\label{eq:1.142}
\end{equation}
注意我们用到了 $R_M$ 的反对称性来改变符号. 观察到 \eqref{eq:1.142} 的右端与坐标的
选取无关. 把 \eqref{eq:1.142} 代入 \eqref{eq:1.132}, 积分并利用 $\Sigma$ 的极小性与 Stokes
定理, 便得到
\begin{equation}
\begin{aligned}
\frac{d^{2}}{dt^{2}}\Big|_{t=0}\Vol\bigl(F(\Sigma,t)\bigr)
&=-\int_{\Sigma}\bigl|\bigl\langle A(\cdot,\cdot),F_t\bigr\rangle\bigr|^{2}
 +\int_{\Sigma}\bigl|\nabla_{\Sigma}^{N}F_t\bigr|^{2}\\
&\quad-\int_{\Sigma}\operatorname{Tr}_{\Sigma}
 \bigl\langle \mathrm{R}_M(\cdot,F_t)\cdot,F_t\bigr\rangle\\
&=-\int_{\Sigma}\bigl\langle F_t,LF_t\bigr\rangle.
\end{aligned}
\label{eq:1.143}
\end{equation}
\end{proof}

这个自伴算子 $L$ 就是所谓的\emph{稳定算子}(或\emph{Jacobi 算子}), 它在
$\Sigma$ 的法向量场 $X$ 上由
\begin{equation}
LX=\Delta_{\Sigma}^{N}X
 +\operatorname{Tr}\bigl[\mathrm{R}_M(\cdot,X)\cdot\bigr]
 +\tilde A(X)
\label{eq:1.144}
\end{equation}
定义, 其中 $\tilde A$ 是由
\begin{equation}
\tilde A(X)=\sum_{i,j=1}^{k}g\bigl(A(E_i,E_j),X\bigr)A(E_i,E_j)
\label{eq:1.145}
\end{equation}
定义的 \emph{Simons 算子}, 而 $\Delta_{\Sigma}^{N}$ 是\emph{法丛上的
Laplace 算子}, 即
\begin{equation}
\Delta_{\Sigma}^{N}X
 =\sum_{i=1}^{k}\bigl(\nabla_{E_i}\nabla_{E_i}X\bigr)^{N}
 -\sum_{i=1}^{k}\bigl(\nabla_{(\nabla_{E_i}E_i)^{T}}X\bigr)^{N}.
\label{eq:1.146}
\end{equation}
满足 $LX=0$ 的法向量场 $X$ 称为 \emph{Jacobi 场}.

若法丛存在整体的正交基, 则称 $\Sigma$ 具有\emph{平凡法丛}. 当 $\Sigma$ 是
超曲面时, 这等价于说 $\Sigma$ 是双侧的.

对具有平凡法丛的超曲面, 稳定算子大为简化, 因为此时它成为作用在函数上的
算子. 即: 若把法向量场 $X=\eta N$ 与 $\eta$ 等同, 则
\begin{equation}
L\eta=\Delta_{\Sigma}\eta+|A|^{2}\eta
 +\operatorname{Ric}_M(N,N)\eta,
\label{eq:1.147}
\end{equation}
其中 $\operatorname{Ric}_M$ 是 $M$ 的 Ricci 张量(例如见 \cite{dC2} 第
97 页).

我们采用如下约定: 若存在不恒为零的法向量场 $X$, 它在 $\partial\Omega$ 上
为零且满足
\begin{equation}
LX+\lambda X=0,
\label{eq:1.148}
\end{equation}
则称 $\lambda$ 是 $L$ 在 $\Omega\subset\Sigma$ 上的一个(Dirichlet)特征值.

\begin{definition}
\label{def:1.31}
紧极小曲面 $\Sigma$ 的 \emph{Morse 指标}是稳定算子 $L$ 在法丛的、在边界上
为零的光滑截面空间上作用的负特征值个数(按重数计).
\end{definition}

第二变分公式表明, 若 $\Sigma^k\subset M^n$ 是极小子流形, 则面积泛函在
$\Sigma$ 处的 Hessian 由
\begin{equation}
-\int_{\Sigma}\langle\cdot,L\,\cdot\rangle
\label{eq:1.149}
\end{equation}
给出. 于是, 我们本可以把 $\Sigma$ 的 Morse 指标等价地定义为 $\Sigma$ 作为
面积泛函临界点的指标.

\mtpara{稳定性}
若对一切边界固定的变分 $F$ 都有
\begin{equation}
\frac{d^{2}}{dt^{2}}\Big|_{t=0}\Vol\bigl(F(\Sigma,t)\bigr)
 =-\int_{\Sigma}\langle F_t,LF_t\rangle\geq0,
\label{eq:1.150}
\end{equation}
则称极小子流形 $\Sigma^k\subset M^n$ 是\emph{稳定的}. 注意, 稳定性等价于
要求稳定算子是负半定的(即 Morse 指标为零). 又注意, 若
$\Sigma^{n-1}\subset\R^n$ 是满足极小曲面方程的函数的图, 则 $\Sigma$ 是稳定
的, 因为此时 $\Sigma$ 事实上是面积最小的. 若一个完全(可以非紧)的无边界极小
子流形的一切紧子区域都稳定, 则称它是稳定的.

对稳定极小超曲面, 我们有下述有用的不等式:

\begin{lemma}[稳定性不等式]
\label{lem:1.32}
设 $\Sigma^{n-1}\subset M^n$ 是具有平凡法丛的稳定极小超曲面, 则对一切具紧
支集的 Lipschitz 函数 $\eta$,
\begin{equation}
\int_{\Sigma}\Bigl(\inf_M\operatorname{Ric}_M+|A|^{2}\Bigr)\eta^{2}
 \leq\int_{\Sigma}\bigl|\nabla_{\Sigma}\eta\bigr|^{2}.
\label{eq:1.151}
\end{equation}
\end{lemma}

\begin{proof}
由于 $\Sigma$ 稳定,
\begin{equation}
0\leq-\int_{\Sigma}\eta L\eta
 =-\int_{\Sigma}\Bigl(\eta\Delta_{\Sigma}\eta+|A|^{2}\eta^{2}
 +\operatorname{Ric}_M(N,N)\eta^{2}\Bigr).
\label{eq:1.152}
\end{equation}
分部积分给出
\begin{equation}
\int_{\Sigma}\Bigl(\operatorname{Ric}_M(N,N)+|A|^{2}\Bigr)\eta^{2}
 \leq\int_{\Sigma}\bigl|\nabla_{\Sigma}\eta\bigr|^{2}.
\label{eq:1.153}
\end{equation}
这就证明了引理.
\end{proof}

当 $M$ 的曲率具有某种正性时, 稳定性不等式给出稳定极小超曲面所受的限制.
下一个推论记下这方面的两个版本(第一个属于 J. Simons, 第二个属于 Schoen 与
Yau):

\begin{corollary}
\label{cor:1.33}
设 $\Sigma^{n-1}\subset M^n$ 是具平凡法丛的闭稳定极小超曲面.
\begin{itemize}
  \item 若 $\operatorname{Ric}_M\geq0$, 则 $\Sigma$ 全测地, 且
    $\operatorname{Ric}_M(N,N)=0$ 于 $\Sigma$.
  \item 若 $\operatorname{Scal}_M>0$ 且 $n=3$, 则 $\Sigma$ 是 $\bS^2$ 或
    $\mathbb{RP}^2$, 且
\begin{equation}
\int_{\Sigma}\Bigl(\operatorname{Scal}_M+|A|^{2}\Bigr)\leq8\pi.
\label{eq:1.154}
\end{equation}
\end{itemize}
\end{corollary}

\begin{proof}
由于 $\Sigma$ 紧且无边界, 我们可以在稳定性不等式中取常函数 $\eta=1$, 得到
\begin{equation}
\int_{\Sigma}\Bigl(\operatorname{Ric}_M(N,N)+|A|^{2}\Bigr)\leq0.
\label{eq:1.155}
\end{equation}
第一个结论立即得出. 对第二个结论, 设
\[
E_1,\ E_2,\ \text{以及}\ E_3=N
\]
是沿 $\Sigma$ 的正交基. 由于 $R_{3333}=0$, 法向 Ricci 曲率为
\begin{equation}
\operatorname{Ric}_M(N,N)=R_{1313}+R_{2323}.
\label{eq:1.156}
\end{equation}
类似地, $\operatorname{Ric}_M(E_1,E_1)=R_{1212}+R_{1313}$,
$\operatorname{Ric}_M(E_2,E_2)=R_{1212}+R_{2323}$. 于是标量曲率为
\begin{equation}
\begin{aligned}
\operatorname{Scal}_M
&=\operatorname{Ric}_M(E_1,E_1)+\operatorname{Ric}_M(E_2,E_2)
 +\operatorname{Ric}_M(N,N)
\label{eq:1.157}
\end{aligned}
\end{equation}
\begin{equation}
\begin{aligned}
&=2\bigl(R_{1212}+R_{1313}+R_{2323}\bigr)
 =2\operatorname{Ric}_M(N,N)+2R_{1212}.
\label{eq:1.158}
\end{aligned}
\end{equation}
Gauss 方程与极小性给出 $K_{\Sigma}=R_{1212}+\det(A)=R_{1212}-|A|^{2}/2$,
故
\begin{equation}
\operatorname{Ric}_M(N,N)=\tfrac12\operatorname{Scal}_M-K_{\Sigma}
 -\tfrac12|A|^{2}.
\label{eq:1.159}
\end{equation}
把它代入 \eqref{eq:1.155}, 便得到
\begin{equation}
\tfrac12\int_{\Sigma}
\bigl(\operatorname{Scal}_M+|A|^{2}\bigr)
 \leq\int_{\Sigma}K_{\Sigma}=2\pi\chi(\Sigma),
\label{eq:1.160}
\end{equation}
其中 $\chi(\Sigma)$ 是 $\Sigma$ 的 Euler 示性数, 等式用了 Gauss--Bonnet
公式(见 \cite{dC1} 第 274 页). 于是 $\chi(\Sigma)>0$, 结论成立.
\end{proof}

利用引理~\ref{lem:1.32} 与对数截断论证(如同定理~\ref{thm:1.21} 的证明),
容易给出定理~\ref{thm:1.21} 的第二个证明. 我们将在下一章回到这一观点.

\mtpara{稳定性的一个刻画}
本节最后, 我们给出具有平凡法丛的极小超曲面稳定性的几个有用刻画, 并导出
一些推论. 这需要比本讲义其余部分更多的 PDE 背景; 遇到这种情况时, 我们会
给出确切的参考文献.

对具有平凡法丛的极小超曲面, 我们已经看到稳定性等价于: 对每个
$\Omega\subset\Sigma$ 都有 $\lambda_1(\Omega,L)\geq0$, 其中
\begin{equation}
\lambda_1(\Omega,L)=\inf\Bigl\{-\int\eta L\eta
 \Bigm|\ \eta\in C_0^{\infty}(\Omega)\ \text{且}\ \int_{\Omega}\eta^{2}=1
 \Bigr\}.
\label{eq:1.161}
\end{equation}
对光滑函数 $u$, 定义 $W^{1,2}$-范数为
\begin{equation}
\|u\|_{W^{1,2}}^{2}=\int u^{2}+\int|\nabla u|^{2}.
\label{eq:1.162}
\end{equation}
Sobolev 空间 $W_0^{1,2}(\Omega)$ 是 $\Omega$ 上具紧支集的光滑函数空间关于
$W^{1,2}$-范数的闭包. 类似地, $W^{1,2}(\Omega)$ 是 $\Omega$ 上光滑函数空间
关于 $W^{1,2}$-范数的闭包.

由标准的椭圆理论(例如见 \cite{HaLi} 或 \cite{GiTr}), 我们得到:

\begin{lemma}
\label{lem:1.34}
若 $L$ 与 $\Omega\subset\Sigma$ 如上, 且 $\lambda_1=\lambda_1(\Omega,L)$, 则
\begin{equation}
\lambda_1=\inf\Bigl\{
 \frac{\int|\nabla_{\Sigma}\eta|^{2}-|A|^{2}\eta^{2}
  -\operatorname{Ric}_M(N,N)\eta^{2}}{\int\eta^{2}}
 \Bigm|\ \eta\in W_0^{1,2}(\Omega)\Bigr\}.
\label{eq:1.163}
\end{equation}
此外, 若 $u\in W_0^{1,2}(\Omega)$ 满足
\begin{equation}
\frac{\int|\nabla_{\Sigma}u|^{2}-|A|^{2}u^{2}
 -\operatorname{Ric}_M(N,N)u^{2}}{\int u^{2}}=\lambda_1,
\label{eq:1.164}
\end{equation}
则 $u$ 光滑且 $Lu=-\lambda_1u$.
\end{lemma}

\begin{proof}
令 $I$ 等于 \eqref{eq:1.163} 的右端. 为记号简洁, 在 $\Sigma$ 上定义函数
\begin{equation}
V(x)=|A|^{2}(x)+\operatorname{Ric}_M(N,N)(x).
\label{eq:1.165}
\end{equation}
由于当 $\eta$ 是具紧支集的 $C^{2}$ 函数时, 散度定理给出
$\int|\nabla_{\Sigma}\eta|^{2}=-\int\eta\Delta_{\Sigma}\eta$, 我们总有
\begin{equation}
I\leq\lambda_1.
\label{eq:1.166}
\end{equation}
另一方面, 取极小化序列 $\eta_j\in W_0^{1,2}(\Omega)$ 使
\begin{equation}
\frac{\int|\nabla_{\Sigma}\eta_j|^{2}-V\eta_j^{2}}{\int\eta_j^{2}}
 \leq I+\frac1j.
\label{eq:1.167}
\end{equation}
把 $\eta_j$ 乘上一个常数后, 可以假定 $\int\eta_j^{2}=1$. 由 Rellich 紧性
定理(\cite{GiTr} 的定理 7.22), 可以取子列(仍记为 $\eta_j$), 使得:
\begin{itemize}
  \item $\eta_j$ 在 $W^{1,2}$ 中弱收敛到某个函数 $\eta\in W_0^{1,2}$;
  \item $\eta_j$ 在 $L^{2}$ 中强收敛到 $\eta$, 特别地 $\int\eta^{2}=1$.
\end{itemize}
此外, 能量在弱收敛下是下半连续的, 故
\[
\int|\nabla_{\Sigma}\eta|^{2}
 \leq\liminf\int|\nabla_{\Sigma}\eta_j|^{2}.
\]
利用 $V$ 是连续函数, 便知 $\eta$ 满足
\begin{equation}
\frac{\int|\nabla_{\Sigma}\eta|^{2}-V\eta^{2}}{\int\eta^{2}}\leq I.
\label{eq:1.168}
\end{equation}
由 $I$ 的定义, \eqref{eq:1.168} 中必须取等号, 特别地, $\eta$ 取到该下确界 $I$.

下面我们证明: $W_0^{1,2}$ 中任何使 \eqref{eq:1.168} 取等号的 $\eta$ 都光滑且满足
$L\eta=-I\eta$. 两个结论都将由此推出. 设 $\eta\in W_0^{1,2}$ 使 \eqref{eq:1.168} 取
等号, 并规范化使
\[
\int\eta^{2}=1.
\]
若 $\psi$ 是 $\Omega$ 中具紧支集的光滑函数, 且
\begin{equation}
\int\psi\eta=0,
\label{eq:1.169}
\end{equation}
则显然
\begin{equation}
\frac{d}{dt}\Big|_{t=0}\int(\eta+t\psi)^{2}=0.
\label{eq:1.170}
\end{equation}
由 \eqref{eq:1.170} 与 $I$ 的定义, 必须有
\begin{equation}
\begin{aligned}
0&=\frac{d}{dt}\Big|_{t=0}\int|\nabla_{\Sigma}(\eta+t\psi)|^{2}
 -V(\eta+t\psi)^{2}\\
&=2\int\bigl(\langle\nabla_{\Sigma}\eta,\nabla_{\Sigma}\psi\rangle
 -V\eta\psi\bigr).
\end{aligned}
\label{eq:1.171}
\end{equation}
由逼近, 方程 \eqref{eq:1.171} 对任何满足 \eqref{eq:1.169} 的 $\psi\in W_0^{1,2}(\Omega)$
成立. 特别地, 给定任意 $\phi\in W_0^{1,2}(\Omega)$, \eqref{eq:1.171} 对
\begin{equation}
\psi=\phi-\eta\int(\phi\eta)
\label{eq:1.172}
\end{equation}
成立, 从而
\begin{equation}
\int\bigl(\langle\nabla_{\Sigma}\eta,\nabla_{\Sigma}\phi\rangle
 -V\eta\phi\bigr)
 =\Bigl(\int\phi\eta\Bigr)\int\bigl(|\nabla_{\Sigma}\eta|^{2}-V\eta^{2}\bigr)
 =I\int(\phi\eta),
\label{eq:1.173}
\end{equation}
其中最后一个等式用了 $\eta$ 使 \eqref{eq:1.168} 取等号以及 $\int\eta^{2}=1$. 由于
\eqref{eq:1.173} 对一切 $\phi\in W_0^{1,2}(\Omega)$ 成立, $\eta$ 是
\[
L\eta=(\Delta_{\Sigma}+V)\eta=-I\eta
\]
的弱解. 于是由椭圆正则性(\cite{GiTr} 的定理 8.8), $\eta$ 必光滑且满足
$L\eta=-I\eta$. 由于 $\eta$ 光滑, 它是 $\lambda_1$ 定义中合法的试验函数,
故
\begin{equation}
\lambda_1\leq I.
\label{eq:1.174}
\end{equation}
把它与(容易得到的)反向不等式结合, 得 $\lambda_1=I$, 即得引理的第一个结论.
既然我们已经证明 $\eta$ 光滑且 $L\eta=-I\eta$, 第二个结论也随之成立.
\end{proof}

把引理~\ref{lem:1.34} 与 Harnack 不等式结合起来, 我们在下一个引理中看到,
第一特征值的任何特征函数都不能变号.

\begin{lemma}
\label{lem:1.35}
若 $u$ 是 $\Omega$ 上在 $\partial\Omega$ 处为零的光滑函数, 且
$Lu=-\lambda_1u$, 其中 $\lambda_1=\lambda_1(\Omega,L)$, 则 $u$ 在 $\Omega$
中不能变号.
\end{lemma}

\begin{proof}
不妨设 $u$ 不恒为零. 由于 $u$ 在 $\partial\Omega$ 上为零, $|u|$ 也是. 事实
上, 容易看出 $|u|$ 同样使 \eqref{eq:1.164} 取到最小值, 从而由
引理~\ref{lem:1.34}, $|u|$ 光滑且 $L|u|=-\lambda_1|u|$. 由于 $|u|\geq0$ 且
$|u|$ 不恒为零, Harnack 不等式(引理~\ref{lem:1.46})蕴含 $|u|>0$ 于
$\Omega$, 引理得证.
\end{proof}

由于不同特征值对应的特征函数彼此正交, 引理~\ref{lem:1.35} 蕴含只有最低的
特征函数不变号, 而且第一特征值的重数为 $1$. 作为推论, 我们看到: 若
$\Sigma\subset M$ 是无边界、具平凡法丛的稳定极小超曲面, 则
$\tilde\Sigma\subset\tilde M$ 也稳定, 其中
\begin{equation}
G:\tilde M\to M
\label{eq:1.175}
\end{equation}
是覆盖映射, $\tilde\Sigma=G^{-1}(\Sigma)$, 而 $\tilde M$ 赋予拉回度量.
另一方面, 简单的例子表明, 稳定极小子流形的覆盖一般并不稳定(例如考虑
$\mathbb{RP}^2\subset\mathbb{RP}^3$).

更一般地, 我们有 Barta 定理 \cite{Ba} 的下述版本:

\begin{lemma}
\label{lem:1.36}
设 $\Sigma$ 是具平凡法丛的极小超曲面, $L$ 是其稳定算子,
$\Omega\subset\Sigma$ 是有界区域. 若 $\Omega$ 上存在正函数 $u$ 使 $Lu=0$,
则 $\Omega$ 稳定.
\end{lemma}

\begin{proof}
令 $q=|A|^{2}+\operatorname{Ric}_M(N,N)$, 使 $L=\Delta_{\Sigma}+q$. 由于
$u>0$, $w=\log u$ 有意义, 且满足
\begin{equation}
\Delta_{\Sigma}w=-q-|\nabla_{\Sigma}w|^{2}.
\label{eq:1.176}
\end{equation}
设 $f$ 是 $\Omega$ 上具紧支集的光滑函数. 把 \eqref{eq:1.176} 两边乘以 $f^{2}$ 并
分部积分, 得
\begin{equation}
\begin{aligned}
\int f^{2}q+\int f^{2}|\nabla_{\Sigma}w|^{2}
&=-\int f^{2}\Delta_{\Sigma}w
 \leq2\int|f|\,|\nabla_{\Sigma}f|\,|\nabla_{\Sigma}w|\\
&\leq\int f^{2}|\nabla_{\Sigma}w|^{2}+\int|\nabla_{\Sigma}f|^{2},
\end{aligned}
\label{eq:1.177}
\end{equation}
其中第二个不等式由 Cauchy--Schwarz 不等式得到. 从 \eqref{eq:1.177} 两边消去
$\int f^{2}|\nabla_{\Sigma}w|^{2}$ 项, 便看出
\begin{equation}
-\int fLf\geq0.
\label{eq:1.178}
\end{equation}
由于这对任何这样的 $f$ 都成立, 引理得证.
\end{proof}

注意, 若 $\Sigma$ 是闭的, 或者更一般地, 若 $u$ 在 $\partial\Sigma$ 上为零,
则引理~\ref{lem:1.36} 立即由引理~\ref{lem:1.35} 推出.

由
\begin{equation}
F(\cdot,t):(x_1,x_2,x_3)\mapsto(x_1,x_2,x_3+t)
\label{eq:1.179}
\end{equation}
给出的变分 $F:\R^3\times\R\to\R^3$ 是 $\R^3$ 的单参数等距群, 从而对任何
曲面 $\Sigma\subset\R^3$, $\Area F(\Sigma,t)$ 是常数. 变分向量场为
$F_t=(0,0,1)$. 若 $\Sigma$ 极小, 则第二变分公式蕴含
$\langle N,(0,0,1)\rangle$(变分向量场的法向分量)是一个 Jacobi 场. 此外,
当 $\Sigma=\Graph_u$ 时,
\begin{equation}
\bigl\langle N,(0,0,1)\bigr\rangle
 =\frac{1}{\sqrt{1+|\nabla_{\R^2}u|^{2}}},
\label{eq:1.180}
\end{equation}
因而是正的 Jacobi 场. 于是, 引理~\ref{lem:1.36} 给出看出极小图稳定性的
另一条途径.

若流形 $\Sigma$ 上任何正超调和函数 $u$(即 $\Delta_{\Sigma}u\leq0$)都是常数,
则称 $\Sigma$ 是\emph{抛物的}. 下一个命题用对数截断技巧说明二次体积增长
蕴含抛物性. 在这个命题中, 我们用 $B_s^{\Sigma}=B_s^{\Sigma}(p)$ 表示
$\Sigma$ 中的内蕴(测地)球.

\begin{proposition}
\label{prop:1.37}
若 $\Sigma$ 是完全曲面, 且对一切 $s>0$ 有
\begin{equation}
\Area\bigl(B_s^{\Sigma}\bigr)\leq Cs^{2},
\label{eq:1.181}
\end{equation}
则 $\Sigma$ 是抛物的.
\end{proposition}

\begin{proof}
设 $u>0$ 且 $\Delta_{\Sigma}u\leq0$, 令
\begin{equation}
w=\log u,
\label{eq:1.182}
\end{equation}
于是 $|\nabla_{\Sigma}w|^{2}\leq-\Delta_{\Sigma}w$. 设 $r$ 表示到 $p$ 的距离,
定义截断函数 $\eta$ 为
\begin{equation}
\eta=
\begin{cases}
1, & \text{若 } r^{2}\leq R,\\
2-\dfrac{\log r^{2}}{\log R}, & \text{若 } R<r^{2}\leq R^{2},\\
0, & \text{若 } r^{2}>R^{2}.
\end{cases}
\label{eq:1.183}
\end{equation}
由 Stokes 定理与吸收不等式 $\bigl(2ab\leq\frac12a^{2}+2b^{2}\bigr)$, 得到
\begin{equation}
\begin{aligned}
\int\eta^{2}|\nabla_{\Sigma}w|^{2}
&\leq-\int\eta^{2}\Delta_{\Sigma}w
 \leq2\int\eta\,|\nabla_{\Sigma}\eta|\,|\nabla_{\Sigma}w|\\
&\leq\frac12\int\eta^{2}|\nabla_{\Sigma}w|^{2}
 +2\int|\nabla_{\Sigma}\eta|^{2}.
\end{aligned}
\label{eq:1.184}
\end{equation}
代入 $\eta$ 的定义与面积界, 得
\begin{equation}
\begin{aligned}
\int_{B_{\sqrt R}^{\Sigma}}|\nabla_{\Sigma}w|^{2}
&\leq\int\eta^{2}|\nabla_{\Sigma}w|^{2}
 \leq4\int|\nabla_{\Sigma}\eta|^{2}\\
&\leq\frac{16}{(\log R)^{2}}
 \sum_{\ell=\frac12\log R}^{\log R}
 \int_{\bigl(B_{\mathrm{e}^{\ell}}^{\Sigma}
 \setminus B_{\mathrm{e}^{\ell-1}}^{\Sigma}\bigr)}r^{-2}\\
&\leq\frac{16}{(\log R)^{2}}
 \sum_{\ell=\frac12\log R}^{\log R}C\mathrm{e}^{2}
 \leq\frac{8C\mathrm{e}^{2}}{\log R}.
\end{aligned}
\label{eq:1.185}
\end{equation}
令 $R\to\infty$, 便得 $w$ 为常数.
\end{proof}

应用命题~\ref{prop:1.37}, 我们看到整个二维极小图 $\Sigma$ 是抛物的. 这是
因为内蕴距离不小于欧氏距离, 从而面积界 \eqref{eq:1.20} 蕴含极小图具有二次面积增长.
按 \eqref{eq:1.180} 取
\begin{equation}
u=\bigl\langle N,(0,0,1)\bigr\rangle
\label{eq:1.186}
\end{equation}
便得到正的 Jacobi 场. 特别地,
\begin{equation}
\Delta_{\Sigma}u
 =-\bigl(|A|^{2}+\operatorname{Ric}_{\R^3}(N,N)\bigr)u
 =-|A|^{2}u\leq0,
\label{eq:1.187}
\end{equation}
故 $u$ 是正超调和函数. 由命题~\ref{prop:1.37}, $u$ 必为常数, 从而
$\Delta_{\Sigma}u=0$, 于是 $|A|^{2}=0$. 换言之, 任何定义在 $\R^2$ 上的完全
极小图必是平坦的. 这给出 Bernstein 定理(定理~\ref{thm:1.21})的另一个证明.

\begin{remark}
\label{rem:1.38}
由 Osserman 定理, 一切具有有限总曲率的极小曲面都是抛物的. 螺旋面与
Riemann 的经典例子具有无限总曲率, 但也是抛物的.
\end{remark}

下面我们给出 Fischer-Colbrie 与 Schoen 的、关于具平凡法丛的完全非紧极小
超曲面稳定性的一个刻画(参看 \cite{Ba} 与 \cite{Du}). 我们假定边界在非空时
是光滑的.

\begin{proposition}[Fischer-Colbrie 与 Schoen, \cite{FiSc}]
\label{prop:1.39}
若 $\Sigma$ 是具平凡法丛的完全非紧极小超曲面, 则下述条件等价:
\begin{equation}
\lambda_1(\Omega,L)\geq0
 \quad\text{对一切有界区域 } \Omega\subset\Sigma;
\label{eq:1.188}
\end{equation}
\begin{equation}
\lambda_1(\Omega,L)>0
 \quad\text{对一切有界区域 } \Omega\subset\Sigma;
\label{eq:1.189}
\end{equation}
\begin{equation}
\text{存在正函数 } u \text{ 使 } Lu=0.
\label{eq:1.190}
\end{equation}
\end{proposition}

\begin{proof}
由引理~\ref{lem:1.36}, \eqref{eq:1.190} 蕴含 \eqref{eq:1.188}.

显然 \eqref{eq:1.189} 蕴含 \eqref{eq:1.188}. 为看出 \eqref{eq:1.188} 与 \eqref{eq:1.189} 的等价性, 给定任意
有界区域 $\Omega_0\subset\Sigma$, 取一个严格更大的有界区域 $\Omega_1$. 特征值
的变分刻画 \eqref{eq:1.161} 蕴含
\begin{equation}
\lambda_1(\Omega_0,L)\geq\lambda_1(\Omega_1,L)\geq0,
\label{eq:1.191}
\end{equation}
其中第二个不等式由 \eqref{eq:1.188} 得到. 设 $u_0$ 表示 $L$ 在 $\Omega_0$ 上的第一
特征函数, 并在 $\Omega_1$ 上定义 $u_1$ 为
\begin{equation}
u_1(x)=
\begin{cases}
u_0(x), & \text{若 } x\in\Omega_0,\\
0, & \text{其他}.
\end{cases}
\label{eq:1.192}
\end{equation}
若 \eqref{eq:1.191} 中取等号, 则由引理~\ref{lem:1.34}, $Lu_1=-\lambda_1u_1$ 于
$\Omega_1$; 又由引理~\ref{lem:1.35}, $u_1>0$ 于 $\Omega_1$. 但这是不可能的,
因为 $u_1$ 在 $\Omega_1\setminus\Omega_0$ 上为零. 于是 \eqref{eq:1.188} 与 \eqref{eq:1.189}
等价.

剩下要证明 \eqref{eq:1.189} 蕴含 \eqref{eq:1.190}. 为此, 固定 $p\in\Sigma$, 对每个 $r>0$ 令
\begin{equation}
B_r^{\Sigma}=B_r^{\Sigma}(p)
 =\bigl\{q\in\Sigma\mid\dist_{\Sigma}(p,q)<r\bigr\}.
\label{eq:1.193}
\end{equation}
由于 $\lambda_1(B_r^{\Sigma},L)>0$, 由 Fredholm 二择一(\cite{GiTr} 的
定理 6.15), 存在唯一的函数 $v_r$ 满足
\begin{equation}
Lv_r=-|A|^{2}-\operatorname{Ric}_M(N,N)\ \ \text{于 } B_r^{\Sigma},
\quad\text{且}\quad v_r=0\ \ \text{于 } \partial B_r^{\Sigma}.
\label{eq:1.194}
\end{equation}
令 $u_r=v_r+1$, \eqref{eq:1.194} 给出
\begin{equation}
Lu_r=0\ \ \text{于 } B_r^{\Sigma},
\quad\text{且}\quad u_r=1\ \ \text{于 } \partial B_r^{\Sigma}.
\label{eq:1.195}
\end{equation}
我们断言
\begin{equation}
u_r>0\ \ \text{于 } B_r^{\Sigma}.
\label{eq:1.196}
\end{equation}
由 Harnack 不等式(引理~\ref{lem:1.46}), 只需证明 $u_r\geq0$ 于
$B_r^{\Sigma}$. 若不然, 我们可以取开集
\begin{equation}
\bigl\{x\in B_r^{\Sigma}\mid u_r(x)<0\bigr\}
\label{eq:1.197}
\end{equation}
的一个非空连通分支 $\Omega$. 按构造, 我们有 $u_r<0$ 于 $\Omega$ 且 $u_r=0$
于 $\partial\Omega$, 而引理~\ref{lem:1.35} 蕴含 $\lambda_1(\Omega,L)=0$.
这个矛盾说明 \eqref{eq:1.196} 成立.

对每个 $r$, 定义正函数
\begin{equation}
w_r=\bigl(u_r(p)\bigr)^{-1}u_r,
\label{eq:1.198}
\end{equation}
并注意 $Lw_r=0$ 且 $w_r(p)=1$.

现在, 设 $K$ 是任一满足 $K\subset B_{R_0}^{\Sigma}$ 的紧集. 应用 Harnack
不等式(直到 $\partial\Sigma$ 的估计见 \cite{GiTr} 的定理 8.27), 对任何
$r\geq2R_0$ 得到
\begin{equation}
\sup_K w_r\leq C_K.
\label{eq:1.199}
\end{equation}
内部与边界的 Schauder 估计(\cite{GiTr} 的定理 6.2 与 6.6)蕴含
\begin{equation}
|w_r|_{C_K^{2,\alpha}}\leq C'_K.
\label{eq:1.200}
\end{equation}
换言之, 若 $K\subset B_{R_0}^{\Sigma}$, 则对一切 $r\geq2R_0$, 各 $w_r$ 具有
一致的 $C_K^{2,\alpha}$ 估计. 由 Arzel\`a--Ascoli 定理, 可以取 $w_r$ 的一个
子列, 使它在紧集上一致地(在 $C^{2,\frac{\alpha}{2}}$ 中)收敛到某个函数 $w$.
这种收敛保证 $w$ 满足 $Lw=0$. 由于每个 $w_r$ 为正且 $w_r(p)=1$, $w$ 非负
且 $w(p)=1$. 最后, Harnack 不等式蕴含 $w$ 也是正的, 这就完成了证明.
\end{proof}

我们可以用命题~\ref{prop:1.39} 给出 Bernstein 定理的一个小幅推广.

\begin{corollary}
\label{cor:1.40}
若 $\Sigma\subset\R^3$ 是无边界的完全、连通、稳定、抛物、双侧极小曲面,
则它必是平面.
\end{corollary}

\begin{proof}
由于 $\Sigma$ 可定向且稳定, 命题~\ref{prop:1.39} 蕴含存在函数 $u>0$ 使
\begin{equation}
\Delta_{\Sigma}u
 =-\bigl(|A|^{2}+\operatorname{Ric}_{\R^3}(N,N)\bigr)u
 =-|A|^{2}u\leq0.
\label{eq:1.201}
\end{equation}
由于 $\Sigma$ 抛物, $u$ 必为常数. 于是 \eqref{eq:1.201} 蕴含 $|A|\equiv0$, 推论得证.
\end{proof}

% ================================================================
\section{多值图}
\label{sec:1.9}

最简单的极小曲面是极小图, 即整张曲面都可以写成平面上某区域上函数的图. 更
一般地, 曲面上每一点都可能有一个邻域是某个固定平面上的图, 但从曲面到该
平面的投影未必是一一的. 此时曲面不是图, 而是一个``多值图''. 对嵌入多值图的
分析在 \cite{CM2}--\cite{CM7} 中起了重要作用.

直观地说, 一个(嵌入的)多值图是这样一张曲面: 在平面上某个圆环的每一点上方,
曲面由 $N$ 个图组成. 为使这一概念精确, 设 $D_r$ 是平面上以原点为心、半径为
$r$ 的圆盘, 并设 $\mathcal{P}$ 是穿孔平面 $\mathbf{C}\setminus\{0\}$ 的
万有覆盖, 具有整体极坐标 $(\rho,\theta)$, 其中 $\rho>0$, $\theta\in\R$.
按照 \cite{CM2}, 圆环 $D_s\setminus D_r$ 上的一个 $N$ 值图, 是函数 $u$ 在
\begin{equation}
\bigl\{(\rho,\theta)\mid r<\rho\leq s,\ |\theta|\leq N\pi\bigr\}
\label{eq:1.202}
\end{equation}
上的单值图.

在实际操作中, 我们通常想象 $\R^3$ 中多叶曲面的直观图像, 并把万有覆盖上的
单值图与它在 $\R^3$ 中的多值像等同起来.

\mtfigure[0.40]{fig-1-24}{多值图. 把两个 $\infty$ 值图沿一条线粘合起来便
  得到螺旋面. 图片提供: Matthias Weber,
  \url{www.indiana.edu/~minimal}.}{fig:1.24}

我们将考虑的多值图都是嵌入的, 这对应于各叶(或各层)之间存在不消失的间隔.
按照 \cite{CM2}, 间隔是指函数(见图~\ref{fig:1.25})
\begin{equation}
w(\rho,\theta)=u(\rho,\theta+2\pi)-u(\rho,\theta).
\label{eq:1.203}
\end{equation}

\mtfigure[0.33]{fig-1-25}{多值极小图的间隔 $w$; 见 \eqref{eq:1.203}.}{fig:1.25}

若 $\Sigma$ 是螺旋面, 则
$\Sigma\setminus\{x_3\text{-轴}\}=\Sigma_1\cup\Sigma_2$, 其中
$\Sigma_1,\Sigma_2$ 是 $\mathbf{C}\setminus\{0\}$ 上的 $\infty$ 值图.
$\Sigma_1$ 是函数 $u_1(\rho,\theta)=\theta$ 的图, $\Sigma_2$ 是函数
$u_2(\rho,\theta)=\theta+\pi$ 的图. ($\Sigma_1$ 是 \eqref{eq:1.46} 中 $s>0$ 的子集,
$\Sigma_2$ 是 $s<0$ 的子集.) 两种情形下间隔都是 $w=2\pi$. 多值极小图是满足
极小曲面方程的函数 $u$ 的多值图.

注意, 对嵌入的多值图, $w$ 的符号决定多值图是按左手方式还是右手方式盘旋,
换言之, 决定向上移动对应于顺时针还是逆时针转动.

容易看出多值图必定稳定. 即: 函数 $\langle N,(0,0,1)\rangle$ 是不消失的
Jacobi 场, 故稳定性立即由引理~\ref{lem:1.36} 推出.

最后, 我们指出间隔 $w$ 满足一个相当好的散度型方程. 即: 引理~\ref{lem:1.26}
蕴含
\begin{equation}
\dive\bigl(a_{i,j}\nabla w\bigr)=0,
\label{eq:1.204}
\end{equation}
其中矩阵 $a_{i,j}$ 有界且正定. 事实上, $a_{i,j}$ 的特征值可以用多值函数
$u$ 的梯度的最大值一致地控制. 注意, 只要 $u$ 的梯度一致有界, 这就足以对
$w$ 得到 Harnack 不等式.

% ================================================================
\section{多值图的局部例子}
\label{sec:1.10}

我们将用 Weierstrass 表示构造一列圆盘
\begin{equation}
\Sigma_i\subset B_1=B_1(0)\subset\R^3,
\label{eq:1.205}
\end{equation}
它们的曲率只在 $0$ 处爆破(见下面定理~\ref{thm:1.41} 的 (1) 与 (2)), 且对
每个 $i$, $\Sigma_i\setminus\{x_3\text{-轴}\}$ 由两个多值图组成; 见 (3).
此外(见 (4)), $\Sigma_i\setminus\{x_3=0\}$ 收敛到两张嵌入极小圆盘
$\Sigma^-\subset\{x_3<0\}$ 与 $\Sigma^+\subset\{x_3>0\}$, 它们各自盘旋进入
$\{x_3=0\}$, 因而都不是固有的; 见图~\ref{fig:1.26}.

\begin{theorem}[Colding--Minicozzi, \cite{CM18}]
\label{thm:1.41}
存在一列紧嵌入极小圆盘 $0\in\Sigma_i\subset B_1\subset\R^3$, 满足
$\partial\Sigma_i\subset\partial B_1$, 且包含竖直线段
$\{(0,0,t)\mid |t|<1\}$, 使得:
\begin{enumerate}
  \item $\lim_{i\to\infty}|A_{\Sigma_i}|^{2}(0)=\infty$;
  \item 对一切 $\delta>0$,
    $\sup_i\sup_{\Sigma_i\setminus B_\delta}|A_{\Sigma_i}|^{2}<\infty$;
  \item 对多值图 $\Sigma_{1,i}$ 与 $\Sigma_{2,i}$, 有
    $\Sigma_i\setminus\{x_3\text{-轴}\}=\Sigma_{1,i}\cup\Sigma_{2,i}$;
  \item $\Sigma_i\setminus\{x_3=0\}$ 收敛到两张嵌入极小圆盘
    $\Sigma^{\pm}\subset\{\pm x_3>0\}$, 满足
    $\overline{\Sigma^{\pm}}\setminus\Sigma^{\pm}
     =B_1\cap\{x_3=0\}$. 此外, 对多值图 $\Sigma_1^{\pm}$ 与
    $\Sigma_2^{\pm}$ 有
    $\Sigma^{\pm}\setminus\{x_3\text{-轴}\}
     =\Sigma_1^{\pm}\cup\Sigma_2^{\pm}$, 其中每个多值图都盘旋进入
    $\{x_3=0\}$; 见图~\ref{fig:1.26}.
\end{enumerate}
\end{theorem}

由 (4) 可知, $\Sigma_i\setminus\{0\}$ 收敛为 $B_1\setminus\{0\}$ 的一个
层状结构(其叶为 $\Sigma^-$、$\Sigma^+$ 以及
$B_1\cap\{x_3=0\}\setminus\{0\}$), 而这个层状结构不能延拓为 $B_1$ 的层状
结构. 即: $0$ 不是可去奇点.

\mtfigure[0.55]{fig-1-26}{定理~\ref{thm:1.41} 中例子的示意图. 去掉
  $x_3$-轴后曲面分成两个多值图; 其中一个用粗线画出.}{fig:1.26}

为证明定理~\ref{thm:1.41}, 我们先构造一族(参数 $a\in(0,1/2)$)极小浸入
$F_a$: 具体选取 \eqref{eq:1.109} 中使用的 Weierstrass 数据
$g=\mathrm{e}^{\ii h_a}$(其中 $h_a=u_a+\ii v_a$)、$\phi=dz$ 与区域
$\Omega_a$. 引理~\ref{lem:1.43} 将证明这族浸入是紧的.
引理~\ref{lem:1.44} 证明各浸入 $F_a:\Omega_a\to\R^3$ 是嵌入.

下一个引理记下 $F$ 的微分.

\begin{lemma}
\label{lem:1.42}
若 $F$ 由 \eqref{eq:1.109} 给出, 其中
$g(z)=\mathrm{e}^{\ii(u(z)+\ii v(z))}$, $\phi=dz$, 则
\begin{equation}
\partial_xF=(\sinh v\cos u,\ \sinh v\sin u,\ 1),
\label{eq:1.206}
\end{equation}
\begin{equation}
\partial_yF=(\cosh v\sin u,\ -\cosh v\cos u,\ 0).
\label{eq:1.207}
\end{equation}
\end{lemma}

对每个 $0<a<1/2$, 令(见图~\ref{fig:1.27})
\begin{equation}
h_a(z)=\frac1a\arctan\Bigl(\frac za\Bigr)
\label{eq:1.208}
\end{equation}
定义在区域
\begin{equation}
\Omega_a=\bigl\{(x,y)\mid |x|\leq1/2,\
 |y|\leq(x^2+a^2)^{3/4}/2\bigr\}
\label{eq:1.209}
\end{equation}
上. 注意 $h_a$ 定义良好, 因为 $\Omega_a$ 单连通且 $\pm\ii a\notin\Omega_a$.
为后文引用, 记
\begin{equation}
\partial_z h_a(z)=\frac{1}{z^2+a^2}
 =\frac{x^2+a^2-y^2-2\ii xy}{(x^2+a^2-y^2)^2+4x^2y^2},
\label{eq:1.210}
\end{equation}
\begin{equation}
\mathrm{K}_a(z)
 =\frac{-|\partial_z h_a|^{2}}{\cosh^{4}v_a}
 =\frac{-|z^2+a^2|^{-2}}
 {\cosh^{4}\bigl(\operatorname{Im}\arctan(z/a)/a\bigr)}.
\label{eq:1.211}
\end{equation}
这里 \eqref{eq:1.211} 用了 \eqref{eq:1.117}. 注意, 由 Cauchy--Riemann 方程,
\begin{equation}
\begin{aligned}
\partial_z h_a
&=(\partial_x-\ii\partial_y)(u_a+\ii v_a)/2
 =\partial_x u_a-\ii\partial_y u_a\\
&=\partial_y v_a+\ii\partial_x v_a.
\end{aligned}
\label{eq:1.212}
\end{equation}

本节余下部分中, 设 $F_a:\Omega_a\to\R^3$ 由 \eqref{eq:1.109} 给出, 其中
$g=\mathrm{e}^{\ii h_a}$, $\phi=dz$, $z_0=0$. 令
$\Omega_0=\cap_a\Omega_a\setminus\{0\}$, 于是
\begin{equation}
\Omega_0=\bigl\{(x,y)\mid 0<|x|\leq1/2,\ |y|\leq|x|^{3/2}/2\bigr\}.
\label{eq:1.213}
\end{equation}
(见图~\ref{fig:1.28}.) 函数族 $h_a$ 不是紧的, 因为对 $z\in\Omega_0$ 有
\[
\lim_{a\to0}|h_a|(z)=\infty.
\]
然而, 下一个引理表明浸入族 $F_a$ 是紧的.

\mtfigurepair{fig-1-27}{区域 $\Omega_a$.}{fig:1.27}
  {fig-1-28}{$\Omega_0=\bigcap_{a>0}\Omega_a\setminus\{0\}$ 及其两个
   分支 $\Omega_0^+$ 与 $\Omega_0^-$.}{fig:1.28}

\begin{lemma}
\label{lem:1.43}
若 $a_j\to0$, 则存在子列 $a_i$, 使 $F_{a_i}$ 在 $\Omega_0$ 的紧子集上
一致地 $C^{2}$ 收敛.
\end{lemma}

\begin{proof}
由于 $h_a$ 与 $-1/z$ 全纯, 且
\begin{equation}
\bigl|\partial_z h_a(z)-\partial_z(-1/z)\bigr|
 =\frac{a^{2}}{|z|^{2}|z^{2}+a^{2}|},
\label{eq:1.214}
\end{equation}
容易得到
\begin{equation}
\lim_{a\to0}\nabla h_a=\nabla(-1/z),
\label{eq:1.215}
\end{equation}
其中收敛在 $\Omega_0$ 的紧子集上是一致的. 由于每个 $v_a(x,0)=0$, 由微积分
基本定理, 各 $v_a$ 在 $\Omega_0$ 的紧子集上一致地 $C^{1}$ 收敛. (遗憾的是,
各 $u_a$ 并不收敛.)

设 $\Omega_0^{\pm}=\{\pm x>0\}\cap\Omega_0$ 是 $\Omega_0$ 的两个分支; 见
图~\ref{fig:1.28}. 令 $b_j^{+}=u_{a_j}(1/2)$,
$b_j^{-}=u_{a_j}(-1/2)$, 并选取子列 $a_i$ 使 $b_i^{-}$ 与 $b_i^{+}$ 都在
模 $2\pi$ 意义下收敛(这是可能的, 因为
$T^{2}=\R^{2}/(2\pi\mathbf{Z}^{2})$ 紧). 照上面的论证,
$h_{a_i}-b_i^{\pm}$ 在 $\Omega_0^{\pm}$ 的紧子集上一致地 $C^{1}$ 收敛.
因此由引理~\ref{lem:1.42}, 对应于 Weierstrass 数据
$g=\mathrm{e}^{\ii(h_{a_i}-b_i^{\pm})}$、$\phi=dz$ 的极小浸入当
$i\to\infty$ 时在 $\Omega_0^{\pm}$ 的紧子集上一致地 $C^{2}$ 收敛.
\end{proof}

证明定理~\ref{thm:1.41} 的主要困难在于证明浸入 $F_a:\Omega_a\to\R^3$ 是
嵌入. 这将容易地从下面的 (A) 与 (B) 推出. 即: 我们在
引理~\ref{lem:1.44} 中证明(见图~\ref{fig:1.29} 与~\ref{fig:1.30}), 对
$|t|\leq1/2$:

(A) 水平切片 $\{x_3=t\}\cap F_a(\Omega_a)$ 是平面中竖直线段 $\{x=t\}$ 的
像, 即 $x_3(F_a(x,y))=x$; 见 (P1).

(B) 像 $F_a(\{x=t\}\cap\Omega_a)$ 是平面 $\{x_3=t\}$ 中某线段上的图(该
线段依赖于 $t$); 见 (P2).

(C) (B) 中图的边界落在球 $B_{r_0}(F_a(t,0))$ 之外, 其中 $r_0>0$, 且对
一切 $a$ 成立; 见 (P3).

\mtfigure[0.62]{fig-1-29}{引理~\ref{lem:1.44} 中的一个水平切片.}{fig:1.29}

\begin{lemma}
\label{lem:1.44}
浸入 $F_a$ 具有下述性质:

(P1) 对每个 $x,y$, 有 $x_3(F_a(x,y))=x$.

(P2) 对每个 $x$, 曲线
\[
F_a(x,\cdot):
\Bigl[-\frac{(x^{2}+a^{2})^{3/4}}{2},
 \frac{(x^{2}+a^{2})^{3/4}}{2}\Bigr]\longrightarrow\{x_3=x\}
\]
是图.

(P3) 存在 $r_0>0$, 使得对每个 $a$ 都有
\[
\Bigl|F_a\Bigl(x,\pm\frac{(x^{2}+a^{2})^{3/4}}{2}\Bigr)
 -F_a(x,0)\Bigr|>r_0.
\]
\end{lemma}

\begin{proof}
由于 $z_0=0$ 且 $\phi=dz$, 由 \eqref{eq:1.109} 得 (P1). 利用 $\Omega_a$ 上
$y^{2}<(x^{2}+a^{2})/4$, \eqref{eq:1.210} 与 \eqref{eq:1.212} 给出
\begin{equation}
\bigl|\partial_y u_a(x,y)\bigr|
 =\frac{2|xy|}{(x^{2}+a^{2}-y^{2})^{2}+4x^{2}y^{2}}
 \leq\frac{4|xy|}{(x^{2}+a^{2})^{2}},
\label{eq:1.216}
\end{equation}
\begin{equation}
\partial_y v_a(x,y)
 =\frac{x^{2}+a^{2}-y^{2}}{(x^{2}+a^{2}-y^{2})^{2}+4x^{2}y^{2}}
 >\frac{3}{8(x^{2}+a^{2})}.
\label{eq:1.217}
\end{equation}
令 $y_{x,a}=(x^{2}+a^{2})^{3/4}/2$. 对 \eqref{eq:1.216} 积分给出
\begin{equation}
\begin{aligned}
\max_{|y|\leq y_{x,a}}\bigl|u_a(x,y)-u_a(x,0)\bigr|
&\leq\int_0^{y_{x,a}}\frac{4|x|t}{(x^{2}+a^{2})^{2}}\,dt\\
&=\frac{|x|}{2(x^{2}+a^{2})^{1/2}}<1.
\end{aligned}
\label{eq:1.218}
\end{equation}
令 $\gamma_{x,a}(y)=F_a(x,y)$. 由于 $v_a(x,0)=0$ 且 $\cos(1)>1/2$, 把
\eqref{eq:1.207} 与 \eqref{eq:1.218} 结合起来给出
\begin{equation}
\begin{aligned}
\bigl\langle\gamma'_{x,a}(y),\gamma'_{x,a}(0)\bigr\rangle
&=\cosh v_a(x,y)\cos\bigl(u_a(x,y)-u_a(x,0)\bigr)\\
&>\cosh v_a(x,y)/2.
\end{aligned}
\label{eq:1.219}
\end{equation}
这里 $\gamma'_{x,a}(y)=\partial_yF_a(x,y)$. 由 \eqref{eq:1.219},
$\gamma'_{x,a}(y)$ 与 $\gamma'_{x,a}(0)$ 的夹角总小于 $\pi/2$; 这给出
(P2).

由于 $v_a(x,0)=0$, 对 \eqref{eq:1.217} 积分给出
\begin{equation}
\min_{y_{x,a}/2\leq|y|\leq y_{x,a}}\bigl|v_a(x,y)\bigr|
 \geq\int_0^{\frac{y_{x,a}}{2}}\frac{3\,dt}{8(x^{2}+a^{2})}
 =\frac{3}{32(x^{2}+a^{2})^{1/4}}.
\label{eq:1.220}
\end{equation}
对 \eqref{eq:1.219} 积分并利用 \eqref{eq:1.220} 给出
\begin{equation}
\bigl\langle\gamma_{x,a}(y_{x,a})-\gamma_{x,a}(0),
 \gamma'_{x,a}(0)\bigr\rangle
 >\frac{(x^{2}+a^{2})^{3/4}}{16}\,
 \mathrm{e}^{(x^{2}+a^{2})^{-1/4}/11}.
\label{eq:1.221}
\end{equation}
由于 $\lim_{s\to0}s^{3}\mathrm{e}^{s^{-1}/11}=\infty$, \eqref{eq:1.221} 及其对
$\gamma_{x,a}(-y_{x,a})$ 的类似式给出 (P3).
\end{proof}

\mtfigure[0.45]{fig-1-30}{引理~\ref{lem:1.44} 中 $F_a(\Omega_a)$ 的水平
  切片.}{fig:1.30}

\begin{corollary}
\label{cor:1.45}
见图~\ref{fig:1.30}. 设 $r_0$ 由 (P3) 给出.
\begin{enumerate}[label=(\roman*)]
  \item $F_a$ 是嵌入.
  \item 当 $|t|<1/2$ 时, $F_a(t,0)=(0,0,t)$.
  \item $\{0<x_1^{2}+x_2^{2}<r_0^{2}\}\cap F_a(\Omega_a)
   =\tilde\Sigma_{1,a}\cup\tilde\Sigma_{2,a}$, 其中
   $\tilde\Sigma_{1,a},\tilde\Sigma_{2,a}$ 是 $D_{r_0}\setminus\{0\}$ 上的
   多值图.
\end{enumerate}
\end{corollary}

\begin{proof}
性质 (P1) 与 (P2) 立即给出 (i).

由于 $z_0=0$, $F(0,0)=(0,0,0)$. 对 \eqref{eq:1.206} 积分并利用 $v_a(x,0)=0$ 便
给出 (ii).

由 \eqref{eq:1.115}, 当 $|g_a|=1$ 时 $F_a$ 是``竖直''的, 即
$\langle N,(0,0,1)\rangle=0$. 然而 $|g_a(x,y)|=1$ 恰当 $y=0$, 从而由
(ii), 像在 $x_3$-轴之外是图. 把它与 (P3) 结合即得 (iii).
\end{proof}

推论~\ref{cor:1.45} 构造出定理~\ref{thm:1.41} 中将要用到的嵌入 $F_a$, 并
证明了性质 (3). 因此, 为证明定理~\ref{thm:1.41}, 我们只需再证明
(1)、(2) 与 (4).

\begin{proof}[定理~\ref{thm:1.41} 的证明]
由伸缩, 只需对某个 $R>0$ 找出一列 $\Sigma_i\subset B_R$.
推论~\ref{cor:1.45} 给出极小嵌入 $F_a:\Omega_a\to\R^3$, 满足当
$|t|<1/2$ 时 $F_a(t,0)=(0,0,t)$, 故对任何 $R\leq r_0$, (3) 成立. 令
\begin{equation}
R=\min\{r_0/2,\ 1/4\},
\label{eq:1.222}
\end{equation}
\begin{equation}
\Sigma_i=B_R\cap F_{a_i}(\Omega_{a_i}),
\label{eq:1.223}
\end{equation}
其中序列 $a_i$ 待定.

为得到 (1), 只需注意由 \eqref{eq:1.211},
\begin{equation}
|\mathrm{K}_a|(0)=a^{-4}\to\infty\quad(a\to0).
\label{eq:1.224}
\end{equation}

下面证明 (2). 首先, 由 \eqref{eq:1.211},
\begin{equation}
\sup_a\sup_{\{|x|\geq\delta\}\cap\Omega_a}
 |\mathrm{K}_a|<\infty
\label{eq:1.225}
\end{equation}
对一切 $\delta>0$ 成立. 结合 (3) 与 Heinz 关于极小图的曲率估计(下一章的
\thmref{thm:2.3}), 便得到 (2).

为得到 (4), 用引理~\ref{lem:1.43} 取 $a_i\to0$, 使映射 $F_{a_i}$ 在紧
子集上一致地 $C^{2}$ 收敛到 $F_0:\Omega_0\to\R^3$. 于是由
引理~\ref{lem:1.44}, $\Sigma_i\setminus\{x_3=0\}$ 收敛到两张嵌入极小圆盘
$\Sigma^{\pm}\subset F_0(\Omega_0^{\pm})$, 且对多值图 $\Sigma_j^{\pm}$ 有
\begin{equation}
\Sigma^{\pm}\setminus\{x_3\text{-轴}\}
 =\Sigma_1^{\pm}\cup\Sigma_2^{\pm}.
\label{eq:1.226}
\end{equation}
为完成证明, 我们证明每个图 $\Sigma_j^{\pm}$ 都是 $\infty$ 值的(从而盘旋
进入 $\{x_3=0\}$). 注意, 由 (3) 与 \eqref{eq:1.207}, 水平集
$\{x_3=x\}\cap\Sigma_j^{\pm}$ 是沿方向
\begin{equation}
\lim_{a\to0}\bigl(\sin u_a(x,0),\ -\cos u_a(x,0),\ 0\bigr)
\label{eq:1.227}
\end{equation}
的直线上的图. 因此, 由于 $0<t<1/4$ 时一个简单的计算给出
\begin{equation}
\lim_{a\to0}\bigl|u_a(t,0)-u_a(2t,0)\bigr|=1/(2t),
\label{eq:1.228}
\end{equation}
我们看到 $\{t<|x_3|<2t\}\cap\Sigma_j^{\pm}$ 含有一个嵌入的 $N_t$ 值图,
其中
\begin{equation}
N_t\approx\frac{1}{4\pi t}\to\infty\quad(t\to0).
\label{eq:1.229}
\end{equation}
于是 $\Sigma_j^{\pm}$ 必盘旋进入 $\{x_3=0\}$, 这就完成了 (4).
\end{proof}

在这个例子之前, 曲率爆破的嵌入极小圆盘序列的局部例子只有伸缩后的螺旋面
序列, 其爆破沿一条线段发生(即沿螺旋面的轴). 现在已有许多相关的例子, 其
爆破集更复杂, 包括 \cite{De2}, \cite{HoWh2}, \cite{Kh} 与 \cite{Kl}.

{\setlength{\emergencystretch}{3em}%
在所有这些局部例子中, 曲率爆破的奇异集都是竖直线的一个子集.
\cite{CM3}--\cite{CM6} 的 Colding--Minicozzi 极限层状结构定理蕴含奇异集
含于一条 Lipschitz 曲线之中(见第 8 章). 在 \cite{MeWe} 中, Meeks 与 Weber
构造了弯曲螺旋面的序列, 其中曲面绕一条弯曲的曲线扭转, 正如螺旋面绕其轴
扭转一样. 取这些曲面当扭数递增时的序列, 便得到一列嵌入极小曲面, 其曲率
沿该弯曲曲线爆破. 这些曲面只在曲线的一个邻域内是嵌入的, 在那里它们构成
双螺旋阶梯. 这种局部的双螺旋阶梯结构其实具有完全的普遍性; 这已在
\cite{CM3}--\cite{CM6} 中证明.\par}

\mtfigurepair{fig-1-31}{有四个扭的 Meeks--Weber 弯曲螺旋面. 图片提供:
  Matthias Weber, \url{www.indiana.edu/~minimal}.}{fig:1.31}
  {fig-1-32}{有七个扭的弯曲螺旋面. 图片提供: Matthias Weber,
   \url{www.indiana.edu/~minimal}.}{fig:1.32}

\mtfigurepair{fig-1-33}{随着扭数增加, 曲率沿一个圆变大. 图片提供:
  Matthias Weber, \url{www.indiana.edu/~minimal}.}{fig:1.33}
  {fig-1-34}{极限将叶状化一个去掉奇异圆的弯曲柱面. 图片提供:
   Matthias Weber, \url{www.indiana.edu/~minimal}.}{fig:1.34}

% ================================================================
\section*{附录: Harnack 不等式}
\addcontentsline{toc}{section}{附录: Harnack 不等式}
\label{sec:1.app-harnack}

下面我们回忆一致椭圆方程非负解的 Harnack 不等式. 我们将使用的是
\cite{GiTr} 定理 8.20 中的版本, 它适用于很一般的一类算子. 对下一个引理,
设 $\mathcal{L}$ 是 $\R^n$ 上由
\begin{equation}
\mathcal{L}u=\sum_{i,j=1}^{n}\frac{\partial}{\partial x_i}
 \Bigl(a^{ij}(x)\frac{\partial u}{\partial x_j}+b^{i}(x)u\Bigr)
 +\sum_{i=1}^{n}c^{i}(x)\frac{\partial u}{\partial x_i}+d(x)u
\label{eq:1.230}
\end{equation}
给出的二阶线性微分算子, 其中系数 $a^{ij},b^{i},c^{i},d$ 是可测函数.

\begin{lemma}
\label{lem:1.46}
设 $\mathcal{L}$ 是 $\Omega\subset\R^n$ 上的二阶线性微分算子, 其有界可测
系数 $a^{ij},b^{i},c^{i},d$ 如 \eqref{eq:1.230}, 且满足
\begin{equation}
\sum_{i,j=1}^{n}a^{ij}x_ix_j\geq\lambda|x|^{2}
\label{eq:1.231}
\end{equation}
其中 $\lambda>0$, 以及
\begin{equation}
\sum_{i,j=1}^{n}\bigl(a^{ij}\bigr)^{2}\leq\Lambda,
\label{eq:1.232}
\end{equation}
\begin{equation}
\lambda^{-2}\sum_{i=1}^{n}
 \bigl(|b^{i}(x)|^{2}+|c^{i}(x)|^{2}\bigr)+\lambda^{-1}|d(x)|\leq\nu^{2},
\label{eq:1.233}
\end{equation}
其中 $\Lambda,\nu<\infty$. 设 $u\in C^{0}(\Omega)\cap W^{1,2}(\Omega)$
在 $\Omega$ 中满足 $u\geq0$, 且 $\mathcal{L}u=0$ 在弱意义下成立. 则对任意
球 $B_{4R}(y)\subset\Omega$, 有
\begin{equation}
\sup_{B_R(y)}u\leq C\inf_{B_R(y)}u,
\label{eq:1.234}
\end{equation}
其中 $C=C\bigl(n,\frac{\Lambda}{\lambda},\nu R\bigr)<\infty$.
\end{lemma}

利用局部坐标与用球链的覆盖论证, 我们可以把引理~\ref{lem:1.46} 的 Harnack
不等式推广到 Riemann 流形中有界区域上的椭圆方程.

% ================================================================
\section*{附录: Bochner 公式}
\addcontentsline{toc}{section}{附录: Bochner 公式}
\label{sec:1.app-bochner}

在 Euclid 空间 $\R^n$ 上, 偏导数可以交换, 因此
\begin{equation}
\frac12\Delta|\nabla u|^{2}
 =|\operatorname{Hess}_u|^{2}
 +\langle\nabla u,\nabla\Delta u\rangle.
\label{eq:1.235}
\end{equation}
在 Riemann 流形 $M$ 上导数不可交换, \eqref{eq:1.235} 一般不再成立. 不过, 导数交换
的失败程度由 Riemann 曲率张量刻画. 利用这一点, Bochner 证明了下面这个极其
有用的公式:

\begin{proposition}[Bochner 公式]
\label{prop:1.47}
在 Riemann 流形 $M$ 上, 有
\begin{equation}
\frac12\Delta|\nabla u|^{2}
 =|\operatorname{Hess}_u|^{2}
 +\langle\nabla u,\nabla\Delta u\rangle
 +\operatorname{Ric}_M(\nabla u,\nabla u).
\label{eq:1.236}
\end{equation}
\end{proposition}

在证明中, 我们将用到 Hessian $\operatorname{Hess}_u$ 的对称性.
Hessian 是把梯度 $\nabla u$ 取共变导数所得的 $2$-张量. 给定向量场 $V$ 与
$W$, 有
\begin{equation}
\operatorname{Hess}_u(V,W)=(\nabla\nabla u)(V,W)
 =\langle\nabla_V\nabla u,W\rangle.
\label{eq:1.237}
\end{equation}
利用联络的度量相容性与梯度的定义, 有
\begin{equation}
V(W(u))=V\langle\nabla u,W\rangle
 =\langle\nabla_V\nabla u,W\rangle
 +\langle\nabla u,\nabla_VW\rangle,
\label{eq:1.238}
\end{equation}
于是 $V$ 与 $W$ 的换位子 $[V,W]$ 作用在 $u$ 上给出
\begin{equation}
\begin{aligned}
[V,W]u&=V(W(u))-W(V(u))\\
&=\operatorname{Hess}_u(V,W)+\langle\nabla u,\nabla_VW\rangle
 -\operatorname{Hess}_u(W,V)-\langle\nabla u,\nabla_WV\rangle\\
&=\operatorname{Hess}_u(V,W)-\operatorname{Hess}_u(W,V)+[V,W]u,
\end{aligned}
\label{eq:1.239}
\end{equation}
其中最后一个等式用了联络的对称性(即用了
$[V,W]=\nabla_VW-\nabla_WV$).

\begin{proof}[命题~\ref{prop:1.47} 的证明]
设 $e_i$ 是 $M$ 的局部正交标架, 于是
\begin{equation}
\Delta v=\sum_i\nabla^{2}v(e_i,e_i)
 =\sum_i\langle\nabla_{e_i}\nabla v,e_i\rangle.
\label{eq:1.240}
\end{equation}
令 $v=\frac12|\nabla u|^{2}$. 我们有
\[
\langle\nabla v,e_j\rangle
 =\bigl\langle\nabla_{e_j}\nabla u,\nabla u\bigr\rangle
 =\operatorname{Hess}_u(e_j,\nabla u)
 =\operatorname{Hess}_u(\nabla u,e_j)
 =\bigl\langle\nabla_{\nabla u}\nabla u,e_j\bigr\rangle,
\]
其中第二与第四个等式是 $\operatorname{Hess}_u$ 的定义, 第三个等式是
$\operatorname{Hess}_u$ 的对称性. 由于这对任何 $e_j$ 成立, 我们断定
\begin{equation}
\nabla v=\nabla_{\nabla u}\nabla u.
\label{eq:1.241}
\end{equation}
把它代入 \eqref{eq:1.240} 给出
\begin{equation}
\begin{aligned}
\Delta v&=\sum_i
 \bigl\langle\nabla_{e_i}\nabla_{\nabla u}\nabla u,e_i\bigr\rangle\\
&=\sum_i\Bigl\langle\Bigl(\nabla_{\nabla u}\nabla_{e_i}\nabla u
 -\nabla_{[\nabla u,e_i]}\nabla u
 +R_M(\nabla u,e_i)\nabla u\Bigr),e_i\Bigr\rangle,
\end{aligned}
\label{eq:1.242}
\end{equation}
其中 $R_M$ 是 $M$ 的 Riemann 曲率张量, 第二个等式是
$R_M(\nabla u,e_i)\nabla u$ 的定义. 和式中最后一项正是
$\operatorname{Ric}_M(\nabla u,\nabla u)$. 第一项可以改写为
\begin{equation}
\begin{aligned}
\sum_i\bigl\langle\nabla_{\nabla u}\nabla_{e_i}\nabla u,e_i\bigr\rangle
&=\sum_i\nabla u\bigl\langle\nabla_{e_i}\nabla u,e_i\bigr\rangle
 -\sum_i\bigl\langle\nabla_{e_i}\nabla u,\nabla_{\nabla u}e_i\bigr\rangle\\
&=\langle\nabla u,\nabla\Delta u\rangle
 -\sum_i\bigl\langle\nabla_{e_i}\nabla u,\nabla_{\nabla u}e_i\bigr\rangle.
\end{aligned}
\label{eq:1.243}
\end{equation}
利用 $u$ 的 Hessian 的对称性, 再利用联络的对称性, 第二项成为
\begin{equation}
\begin{aligned}
-\sum_i\bigl\langle\nabla_{[\nabla u,e_i]}\nabla u,e_i\bigr\rangle
&=-\sum_i\bigl\langle\nabla_{e_i}\nabla u,[\nabla u,e_i]\bigr\rangle\\
&=\sum_i\bigl\langle\nabla_{e_i}\nabla u,
 \nabla_{e_i}\nabla u-\nabla_{\nabla u}e_i\bigr\rangle\\
&=|\operatorname{Hess}_u|^{2}
 -\sum_i\bigl\langle\nabla_{e_i}\nabla u,\nabla_{\nabla u}e_i\bigr\rangle.
\end{aligned}
\label{eq:1.244}
\end{equation}
把三项的表达式代回 \eqref{eq:1.242}, 给出
\begin{equation}
\begin{aligned}
\Delta v&=|\operatorname{Hess}_u|^{2}
 +\langle\nabla u,\nabla\Delta u\rangle
 +\operatorname{Ric}_M(\nabla u,\nabla u)\\
&\quad-2\sum_i
 \bigl\langle\nabla_{e_i}\nabla u,\nabla_{\nabla u}e_i\bigr\rangle.
\end{aligned}
\label{eq:1.245}
\end{equation}
为完成证明, 我们将证明
\begin{equation}
\sum_i\bigl\langle\nabla_{e_i}\nabla u,\nabla_{\nabla u}e_i\bigr\rangle=0.
\label{eq:1.246}
\end{equation}
证明 \eqref{eq:1.246} 最快的办法是指出: 这个量不依赖于标架的选取(因为 \eqref{eq:1.245} 中
其他量都不依赖), 而在给定点处 $\nabla_{e_j}e_i=0$ 的法坐标下它为零. 由于
这个证明有些神秘, 我们再给出 \eqref{eq:1.246} 的一个直接证明. 为此, 把
$\nabla_{e_i}\nabla u$ 与 $\nabla_{\nabla u}e_i$ 按基展开为
\begin{equation}
\nabla_{e_i}\nabla u=\sum_j a_{ij}e_j,
\label{eq:1.247}
\end{equation}
\begin{equation}
\nabla_{\nabla u}e_i=\sum_j b_{ij}e_j,
\label{eq:1.248}
\end{equation}
于是 \eqref{eq:1.246} 的左端就是 $\sum_{i,j}(a_{ij}b_{ij})$. \eqref{eq:1.246} 的证明现在由
以下三个事实推出:
\begin{enumerate}
  \item $a_{ij}=\langle\nabla_{e_i}\nabla u,e_j\rangle$ 对称, 即
    $a_{ij}=a_{ji}$;
  \item $b_{ij}=\langle\nabla_{\nabla u}e_i,e_j\rangle$ 反对称, 即
    $b_{ij}=-b_{ji}$;
  \item 若 $a_{ij}$ 对称而 $b_{ij}$ 反对称, 则
    $\sum_{i,j}(a_{ij}b_{ij})=0$.
\end{enumerate}
事实 (1) 就是 $u$ 的 Hessian 的对称性. 事实 (2) 由度量相容性推出, 因为
$\nabla u\langle e_i,e_j\rangle=0$. 为证明 (3), 注意
\begin{equation}
\sum_{i,j}(a_{ij}b_{ij})=\sum_{i,j}(a_{ji}b_{ji})
 =-\sum_{i,j}(a_{ij}b_{ij}),
\label{eq:1.249}
\end{equation}
故此式为零, 证明完成.
\end{proof}
